\medterm{Vaccination} Giving a vaccine to train the immune system to recognize a specific infection. Vaccination lowers the risk of illness, severe complications, and sometimes transmission; the recommended schedule depends on age, health, and local guidance.

\medterm{Vaccination Schedule} A time-based plan for administering vaccines according to age, previous doses, medical conditions, exposure risks, and public-health guidance. Schedules vary by country and are updated as recommendations change.

\medterm{Vaccinations/Immunizations During Pregnancy} Vaccines given during pregnancy are selected according to safety, disease risk, season, exposure, and local recommendations. Inactivated vaccines such as influenza and Tdap are commonly used when indicated, while some live vaccines are deferred until after pregnancy.
\medterm{Vaccine} A vaccine trains the immune system to recognize a specific infection or toxin and respond more effectively. Vaccination lowers the risk of severe disease and helps protect communities.

\medterm{Vaccine Allergy} An allergic reaction to a vaccine or one of its ingredients. Reactions may be local and mild or, rarely, involve hives, swelling, breathing difficulty, or anaphylaxis; suspected reactions are assessed to identify the responsible ingredient and guide future vaccination.

\medterm{Vaccine-Associated Paralytic Poliomyelitis} An extremely rare form of polio paralysis caused by a weakened poliovirus used in the oral polio vaccine changing back to a harmful form.

\synonyms
VAPP

\medterm{Vaccines} Preparations that train the immune system to recognize a particular germ or toxin before exposure. They reduce the risk of infection and severe illness and may also reduce complications and spread.

\medterm{Vaccines and Immunizations} Vaccines train the immune system to recognize specific pathogens or toxins and reduce the risk of infection, severe disease, complications, and transmission.

\medterm{Vaccinia} A virus used to make smallpox vaccines. Rare accidental infection can cause a sore at the contact site or more widespread illness, especially in people with weakened immunity.

\medterm{Vaccinia Virus} Vaccinia virus is a poxvirus used to make smallpox vaccines and studied in some experimental cancer treatments; accidental infection is uncommon but possible.

\medterm{VACTERL Association} A pattern of birth defects affecting several of the spine, anus, heart, food pipe, windpipe, kidneys, or limbs. There is no single test; the pattern of differences is used to identify it and rule out other conditions.

\medterm{Vacuum-Assisted Delivery} A vaginal birth assisted by a suction cup placed on the baby’s scalp when labor is not progressing or the baby needs to be born promptly. It is used only when the cervix is fully open and the baby’s position and head size are suitable. Possible complications include scalp injury, bleeding, or tears.

\medterm{Vacuum Curettage} A procedure that removes tissue from the uterus using suction, commonly for uterine evacuation or endometrial sampling. Risks include bleeding, infection, retained tissue, and rarely uterine injury.

\medterm{Vacuum Delivery} Assisted vaginal birth using a suction cup placed on the baby’s scalp. It may be used when labor is not progressing or the baby needs to be born promptly and can cause scalp injury, bleeding, or tears.

\medterm{Vacuum Erection Devices} Devices with a cylinder and pump that draw blood into the penis to produce an erection. A flexible ring may hold the erection during sex; bruising, numbness, or discomfort can occur, especially if the ring is left on too long.

\synonyms
Vacuum Erection Pump; Penis Pump

\medterm{Vacuum Extractor} A cup-shaped device applied to the fetal scalp to assist vaginal delivery when specific maternal or fetal indications are present.

\medterm{Vagal Maneuver} A physical action that affects the vagus nerve and can slow the heart’s electrical activity. Examples include bearing down and, in clinical settings, certain pressure or cold-stimulation methods. These maneuvers may help assess or stop some fast heart rhythms.
\synonyms
Vagal Stimulation Technique; Parasympathetic Provocative Maneuver

\medterm{Vagal Nerve Stimulation} A treatment that uses an implanted device to send small electrical signals to the vagus nerve in the neck. It may be used for selected cases of epilepsy or depression that have not improved with other treatments.

\synonyms
VNS

\medterm{Vagal Paraganglioma} A usually slow-growing tumor arising from nerve-related cells near the vagus nerve in the neck. It may cause a neck mass, hoarseness, swallowing difficulty, or nerve symptoms and can be identified with imaging and specialist assessment.

\medterm{Vagina} A muscular, elastic canal connecting the cervix to the vaginal opening. It allows menstrual blood to leave the body, receives a penis or other objects during vaginal sex, and serves as the birth canal.

\medterm{Vagina Carcinoma} Cancer arising from the vagina. It may cause unusual bleeding, discharge, pain, or a visible sore, although early cancer may cause no symptoms; examination and tissue testing are needed for diagnosis.

\medterm{Vaginal Agenesis} A condition in which the vagina is absent or has not fully formed before birth. It may occur with differences in the uterus, kidneys, hearing, or spine; care is individualized and may include counseling, gradual treatment, or surgery when wanted.

\medterm{Vaginal Atresia} Absence or complete blockage of the vaginal canal from birth. It may cause absent menstrual flow, repeated monthly pelvic pain, or difficulty with penetration.

\medterm{Vaginal Atrophy} Thinning, dryness, and irritation of vaginal tissues, usually caused by lower estrogen after menopause or other hormone changes. It can cause burning, itching, pain during sex, bleeding, or urinary symptoms and may improve with moisturizers, lubricants, or prescribed hormone treatment.

\synonyms
Atrophic Vaginitis; Vulvovaginal Atrophy; Genitourinary Syndrome of Menopause; GSM; Vaginal Dryness

\medterm{Vaginal Bleeding} Unexpected, unusually heavy, or prolonged bleeding from the vagina, including bleeding between periods, after sexual intercourse, during pregnancy, or after menopause. Causes include hormonal changes, infection, pregnancy-related conditions, structural lesions, medications, and cancer.

\medterm{Vaginal Bleeding Between Periods} Bleeding between menstrual periods may result from contraception, infection, pregnancy, polyps, fibroids, hormonal changes, or cancer. Persistent or heavy bleeding can have a serious cause.

\medterm{Vaginal Bleeding During Early Pregnancy} Bleeding from the vagina during the first part of pregnancy can result from cervical changes, miscarriage, ectopic pregnancy, infection, or other conditions. A small amount may occur without a serious cause, but bleeding with heavy loss, severe or one-sided pain, shoulder pain, dizziness, or fainting can indicate a serious problem.

\medterm{Vaginal Bleeding During Late Pregnancy} Bleeding from the vagina later in pregnancy may be caused by placenta previa, placental separation, labor, cervical disease, or other complications. Significant bleeding can be a serious pregnancy complication.

\medterm{Vaginal Bleeding In Early Pregnancy} Alternate name; see \textit{Vaginal Bleeding During Early Pregnancy}.

\medterm{Vaginal Bleeding In Late Pregnancy} Alternate name; see \textit{Vaginal Bleeding During Late Pregnancy}.

\medterm{Vaginal Bleeding In Pregnancy} Broad index label for bleeding during pregnancy; see \textit{Vaginal Bleeding During Early Pregnancy} and \textit{Vaginal Bleeding During Late Pregnancy}. Heavy bleeding or bleeding with pain, dizziness, fainting, or reduced fetal movement can indicate a serious pregnancy complication.

\medterm{Vaginal Cancer} A rare cancer of the vagina, most often arising from its surface lining. It may cause bleeding, discharge, pain, or a lump; treatment depends on its type, size, location, and spread.

\synonyms
Vaginal Carcinoma; Malignant Vaginal Neoplasm

\medterm{Vaginal Candidiasis} Alternate name; see \textit{Yeast infection (vaginal)}.

\medterm{Vaginal Cuff} The upper end of the vagina that is closed with stitches after the uterus is removed during a hysterectomy.

\medterm{Vaginal Cuff Dehiscence} An opening of the stitched upper end of the vagina after hysterectomy. It may cause pelvic pain, bleeding, discharge, or tissue bulging through the opening.

\medterm{Vaginal Cysts} Vaginal cysts are fluid-filled or tissue-lined sacs in or beside the vaginal wall, commonly benign but sometimes causing pain, pressure, infection, or obstruction.

\medterm{Vaginal Dilator} A smooth device used to gently stretch or keep open the vaginal canal when it has become narrow or short because of scarring, surgery, radiation, or a difference present from birth.

\medterm{Vaginal Discharge} Vaginal discharge may be normal or may indicate infection, inflammation, hormonal change, or another condition. Normal discharge is usually clear to whitish and odourless and varies with the menstrual cycle, pregnancy, sexual arousal, and some contraceptives.

\medterm{Vaginal Dryness} Reduced vaginal lubrication, commonly associated with menopause, breastfeeding, hormonal
changes, medications, or irritation. It may cause discomfort, itching, painful
intercourse, or minor bleeding.

\medterm{Vaginal Fistula} An abnormal passage between the vagina and another organ, commonly the bladder or rectum. It may cause continuous urine or stool leakage, infection, skin irritation, and emotional distress.

\medterm{Vaginal Foreign Body} An object retained in the vagina, most often a tampon or other inserted material. It may cause discharge, odor, bleeding, pain, or infection.

\medterm{Vaginal Fornix} One of the recesses surrounding the cervix where the vagina meets the cervix. The posterior fornix is adjacent to the peritoneal cul-de-sac and is relevant in pelvic examination and some procedures.

\medterm{Vaginal Fusion And Duplication} Birth-related differences in which the vaginal canal is partly joined, divided, or duplicated. They may occur with abnormalities of the uterus, cervix, urinary tract, or external genitalia and may cause obstruction, pain, menstrual problems, or difficulty with intercourse.

\medterm{Vaginal Health} Vaginal health includes normal vaginal tissues, fluids, microorganisms, and comfort. Persistent odor, itching, pain, bleeding, or unusual discharge can indicate infection or another condition.

\medterm{Vaginal Hematoma} A collection of blood within the vaginal wall or surrounding tissues, usually caused by childbirth, trauma, or surgery and potentially producing severe pain, swelling, or blood loss.

\medterm{Vaginal Hemorrhage} Heavy or clinically significant bleeding from the vagina. Causes include childbirth, pregnancy complications, trauma, infection, structural disease, hormonal disorders, and gynecologic cancer; severe or pregnancy-related bleeding can be life-threatening.

\medterm{Vaginal Hypoplasia} Vaginal hypoplasia is underdevelopment of the vagina, usually present from birth. It may occur with other reproductive or urinary-tract differences and can affect menstruation, sexual function, or fertility.

\medterm{Vaginal Hysterectomy} Surgical removal of the uterus through the vagina. It may be used for selected benign uterine conditions and usually avoids an abdominal incision, although suitability depends on anatomy and disease.

\medterm{Vaginal Inflammation} Alternate name; see \textit{Vaginitis}.

\medterm{Vaginal Intercourse} The insertion of the penis into the vagina for sexual pleasure or reproduction. The most
common form of sexual intercourse. Vaginal intercourse can lead to pregnancy if sperm
are deposited in the vagina.

\medterm{Vaginal Intraepithelial Neoplasia} Abnormal, precancerous cells in the surface lining of the vagina, often related to human papillomavirus. The changes vary from mild to severe and may require monitoring, treatment, or removal.

\medterm{Vaginal Introitus} The entrance to the vagina. It is part of the external genital area and may be examined when assessing pain, infection, injury, scarring, or pelvic-floor problems.

\medterm{Vaginal Itching} Itching or irritation of the vulva or vagina caused by candidiasis, bacterial vaginosis, sexually transmitted infections, dermatitis, hormonal changes, or other conditions. It may occur with discharge, odor, soreness, or burning.

\medterm{Vaginal Laceration} A tear in the vaginal tissue, often caused by childbirth, sexual activity, trauma, or a medical procedure. It may cause pain, swelling, or bleeding.

\medterm{Vaginal Lubrication} Natural moisture made by the vaginal tissues and nearby glands. It usually increases with sexual arousal and may decrease with hormonal changes, some medicines, illness, or irritation.

\medterm{Vaginal Membrane} A thin fold of tissue associated with the vagina. The term may refer to the hymen at the vaginal opening or to an abnormal membrane that blocks the vaginal canal.

\medterm{Vaginal Microbiome} The community of microorganisms living in the vagina, commonly dominated by Lactobacillus species that help maintain an acidic environment. Changes in this community may contribute to bacterial vaginosis, candidiasis, or other symptoms.

\medterm{Vaginal Mucosa} The moist inner lining of the vagina. It contains folds that allow the vaginal canal to stretch and helps protect the tissue.

\medterm{Vaginal Odor} An unusual or unpleasant odor from vaginal discharge that may result from bacterial vaginosis, infection, retained material, hormonal changes, or other conditions.

\medterm{Vaginal Opening} The external entrance to the vagina. It is part of the vulvar vestibule and may be affected by infection, inflammation, trauma, scarring, or congenital variation.

\medterm{Vaginal Pessary} A removable device placed in the vagina to support the vaginal walls or nearby pelvic organs, usually for pelvic-organ prolapse or some forms of urine leakage.

\medterm{Vaginal pH} A measure of how acidic the vagina is. The vaginal environment is normally mildly acidic, and the pH can change with hormones, semen, blood, medicines, or infection.

\medterm{Vaginal Prolapse} Downward bulging of a vaginal wall because the tissues supporting the vagina have weakened. The bulge may involve the bladder, rectum, or the top of the vagina after hysterectomy.

\medterm{Vaginal Rugae} Soft folds in the inner lining of the vagina that allow it to stretch and return toward its usual shape.

\medterm{Vaginal Seeding} Placing vaginal fluid on a baby born by cesarean section in an attempt to expose the baby to the parent’s vaginal microorganisms. Clear health benefits have not been established, and harmful infections may be transmitted.

\medterm{Vaginal Speculum} An instrument inserted into the vagina and gently opened to separate the vaginal walls and visualize the vagina and cervix. It may be used for pelvic examination, cervical cytology such as a Pap test, swab collection, assessment of bleeding or discharge, and selected procedures. Size and design are chosen according to anatomy, comfort, and the clinical purpose.

\medterm{Vaginal Stenosis} Narrowing or shortening of the vaginal canal, which may result from surgery, radiation, scarring, inflammation, trauma, or congenital development. It can cause pain, difficult examination, or difficulty with penetration.

\medterm{Vaginal Swab} A sample of fluid or cells collected from the vagina for laboratory testing, often to look for infection or other changes.

\medterm{Vaginal Trauma} Injury to the vaginal tissues caused by childbirth, sexual activity, foreign objects, procedures, or blunt or penetrating trauma. It may cause bleeding, pain, swelling, or urinary symptoms; severe bleeding or deep injury can be serious.

\medterm{Vaginal Ultrasound Probe} A specialized endocavitary ultrasound transducer used through the vagina to produce images of the uterus, cervix, ovaries, nearby pelvic structures, or an early pregnancy. It sends and receives sound waves and does not use ionizing radiation.

\synonyms
Endovaginal Transducer
Transvaginal Ultrasound Transducer

\medterm{Vaginal Vault Prolapse} Descent of the top of the vagina into or through the vaginal opening, usually after removal of the uterus. It may cause pelvic pressure, a bulge, backache, or difficulty emptying the bladder or bowel; pelvic-floor therapy, a support device, or surgery may help.

\synonyms
Post-Hysterectomy Vault Prolapse; Apical Prolapse

\medterm{Vaginal Vestibule} The space of the vulva between the labia minora that contains the external openings of the urethra and vagina and the ducts of the greater and lesser vestibular glands.

\medterm{Vaginal Wall} The layers of muscle, connective tissue, and moist lining that form the walls of the vaginal canal.

\medterm{Vaginal Yeast Infection} Alternate name; see \textit{Vulvovaginal Candidiasis}.

\medterm{Vagina Neoplasm} An abnormal growth in the vagina that may be benign, precancerous, or cancerous. It can cause bleeding, discharge, pain, or a lump, although early growths may cause no symptoms.

\medterm{Vaginectomy} Surgical removal of part or all of the vagina, usually to treat selected cancers or severe disease; effects may include changes in sexual function, fertility, urination, and bowel function.

\medterm{Vaginismus} Alternate name; see \textit{Genito-Pelvic Pain/Penetration Disorder}.

\medterm{Vaginitis} Inflammation or irritation of the vagina that may cause discharge, itching, burning, odor, pain, or discomfort with urination or sex. Causes include infections, hormones, chemicals, allergies, and tissue dryness.

\medterm{Vaginoplasty} Vaginoplasty is surgery to create, repair, or reshape the vagina. It may be performed for congenital differences, injury, cancer treatment, pelvic-floor problems, or gender-affirming care.

\medterm{Vaginospasm} Alternate name; see \textit{Genito-Pelvic Pain/Penetration Disorder}.

\medterm{Vagotomy} Surgery that cuts selected branches of the vagus nerve to reduce signals that stimulate stomach acid and movement. It was once used for difficult stomach ulcers but is now uncommon because medicines and less invasive treatments are available.

\medterm{Vagus Nerve} The tenth cranial nerve, carrying signals between the brain and the heart, lungs, throat, voice box, and digestive tract. It helps regulate heart rate, breathing, swallowing, voice, and digestion.

\medterm{Vagus Nerve Stimulation} A treatment that sends regular, mild electrical pulses to the vagus nerve through an implanted device in the chest or, in some settings, a handheld device. It is used for selected epilepsy, treatment-resistant depression, and other conditions; benefits may take time and side effects can include voice change, cough, or throat discomfort.

\medterm{Valacyclovir} An antiviral medicine that the body changes into acyclovir to stop herpes viruses from multiplying. It treats cold sores, genital herpes, and shingles; the dose may need adjustment in kidney disease, and excessive levels can affect the nervous system.

\medterm{Valganciclovir} An antiviral medicine used mainly to prevent or treat cytomegalovirus infection after transplantation. The body changes it into ganciclovir, which blocks viral multiplication; it can lower blood counts, injure the kidneys, affect fertility, and harm a fetus.

\medterm{Valgus} A limb or joint deformity in which the part farther from the body bends outward, away from the body’s center line. Examples include a bunion and knock-knees.

\medterm{Valgus Stress Test} A physical examination manoeuvre used to assess the structural integrity of medial collateral ligaments, most commonly at the knee or elbow. A medial-directed force is applied to the joint while stabilizing the limb; excessive joint opening, laxity, or pain indicates a medial collateral ligament sprain or rupture.

\medterm{Valinemia} A rare inherited condition in which the body cannot process the amino acid valine normally. High valine levels may cause feeding problems, low muscle tone, developmental delay, or other nervous-system symptoms.

\medterm{Valley Fever} The common name for coccidioidomycosis, a fungal infection caused by Coccidioides
species and acquired by inhaling airborne spores from soil. It usually causes a mild
respiratory illness but can disseminate to skin, bone, joints, or meninges, especially
in people at higher risk.

\medterm{Valproate} An antiseizure medicine also used for selected mood disorders. It can affect the liver, pancreas, blood cells, weight, and fetal development, so age, pregnancy possibility, other illnesses, and medicines must be considered.

\synonyms
Valproic Acid; Divalproex Sodium

\medterm{Valproate Toxicity} Overdose or adverse exposure causing CNS depression, vomiting, hyperammonemic encephalopathy, hepatotoxicity, pancreatitis, thrombocytopenia, metabolic acidosis, or cerebral edema. Measure serum valproate, glucose, ammonia, liver tests, and acid-base status; L-carnitine may be considered and severe poisoning may require hemodialysis.

\medterm{Valproic Acid} A medicine used for certain seizures, bipolar disorder, and migraine prevention. It can cause liver injury, pancreatitis, birth defects, and developmental harm during pregnancy.

\medterm{Valproic Acid Toxicity} Acute overdose or adverse metabolic toxicity of the broad-spectrum antiepileptic agent valproate, characterized clinically by central nervous system depression, high anion gap metabolic acidosis, hyperammonemia (due to urea cycle inhibition) with encephalopathy, and drug-induced microvesicular hepatic necrosis; treated with L-carnitine.

\medterm{Valrubicin} A chemotherapy medicine placed directly into the bladder for selected bladder cancers that have not responded adequately to other treatment. It can cause bladder irritation, blood in the urine, infection, and, rarely, effects throughout the body.

\medterm{Valsalva} Straining while trying to breathe out against a closed throat, like bearing down during a bowel movement. It changes pressure in the chest and can affect heart rate and blood pressure, so it is not safe for everyone.

\medterm{Valsalva Maneuver} Alternate name; see \textit{Valsalva manoeuvre}.

\medterm{Valsalva Retinopathy} Preretinal (subhyaloid or sub-internal limiting membrane) hemorrhage occurring within the macula following a sudden, dramatic rise in intrathoracic or intra-abdominal pressure (such as from forceful coughing, vomiting, or heavy lifting) that transmits back through the valveless venous system into retinal capillaries.

\medterm{Valsalva's Manoeuvre} Forceful straining while trying to breathe out against a closed airway, like bearing down during a bowel movement. It changes pressure in the chest and can affect heart rate and blood pressure; it is used in selected tests but may be unsafe for some people.

\synonyms
Valsalva Maneuver

\medterm{Valsartan} Valsartan is an angiotensin-II receptor blocker used for hypertension, heart failure, and selected cardiovascular indications; hyperkalaemia, renal-function changes, hypotension, and fetal toxicity are important risks.

\medterm{Valves} Structures that permit fluid or blood to move predominantly in one direction and prevent backward flow; cardiac valves and venous valves are important examples.

\medterm{Valvotomy} Surgical or catheter-based division of narrowed valve tissue to improve blood flow, most often performed for selected stenotic heart valves.

\medterm{Valvulae Conniventes} Large, permanent transverse circular folds of the mucous membrane and submucosa that project into the lumen of the small intestine (most prominent and numerous in the jejunum), which do not disappear with intestinal distension and act to triple the mucosal surface area for absorption.

\synonyms
Plicae Circulares; Kerckring Folds

\medterm{Valvular Disease} Alternate term; see \textit{Heart valve disease}.

\medterm{Valvular Endocarditis} Alternate name; see \textit{Endocarditis}.

\medterm{Valvular Heart Disease} A broad spectrum of structural and functional disorders affecting one or more of the four cardiac valves (aortic, mitral, tricuspid, pulmonic), classified hemodynamically as valvular stenosis (narrowing and outflow obstruction) or valvular regurgitation (insufficiency and backward flow).

\synonyms
VHD

\medterm{Valvular Insufficiency} Incomplete closure of a heart valve, allowing blood to leak backward when the heart contracts or relaxes. It can affect the mitral, aortic, tricuspid, or pulmonary valve and may eventually strain the heart.

\medterm{Valvuloplasty} A procedure used to widen or repair a narrowed or abnormal heart valve, by balloon catheter or open surgery depending on the valve and condition.

\medterm{Valvulotomy} A procedure that cuts open a narrowed heart valve to improve blood flow. It may be performed with a catheter or surgery; the valve can leak afterward and narrowing may recur.

\medterm{Vancomycin} An antibiotic used for selected serious bacterial infections and, by mouth, for certain infections of the large intestine. It can cause infusion reactions, kidney injury, hearing problems, or low blood counts and requires medical monitoring.

\medterm{Vancomycin Flushing Reaction} A reaction that can occur when vancomycin is given too quickly into a vein. It may cause flushing, warmth, itching, and redness of the face, neck, and upper body.

\synonyms
Red Man Syndrome; Red Neck Syndrome

\medterm{Vancomycin-Resistant Enterococcus} Enterococcus bacteria that are resistant to vancomycin, usually because of altered cell-wall precursors. VRE can
cause healthcare-associated urinary, bloodstream, wound, or other infections and may
require susceptibility-guided treatment.

\medterm{Van der Woude Syndrome} An autosomal dominant congenital disorder caused by mutations in the IRF6 gene, characterized by the combination of paramedian lower-lip pits (fistulae), cleft lip, and cleft palate, representing the most frequent form of syndromic orofacial clefting.
\synonyms
VWS; Lip-Pit Syndrome

\medterm{Vandetanib} A targeted cancer medicine used for selected medullary thyroid cancers. It blocks signals that help cancer cells grow and blood vessels form; diarrhea, high blood pressure, skin effects, and dangerous heart-rhythm changes require monitoring.

\medterm{Vaping} Vaping uses an electronic device to heat a liquid and produce an aerosol for inhalation. Aerosols may contain nicotine and other substances that can harm the lungs and cause dependence.

\medterm{Vaporiser} A device that converts a liquid medicine or volatile substance into a vapor for inhalation or controlled delivery, such as in anesthesia.

\medterm{Variable Decelerations} Brief drops in a baby’s heart rate that vary in timing and shape, often caused by pressure on the umbilical cord during labor. Occasional drops may be harmless, but frequent or prolonged drops require assessment of the baby’s oxygen supply.

\medterm{Variant Angina Pectoris} Chest pain caused by a temporary tightening of an artery supplying the heart, often occurring at rest or at night. Episodes may be severe and can disturb the heart rhythm; sudden or prolonged chest pain needs emergency care.

\synonyms
Prinzmetal Angina; Vasospastic Angina

\medterm{Variant CJD} Alternate name; see \textit{Creutzfeldt-Jakob disease}.

\medterm{Variant Classification} The process of estimating whether a genetic change is harmless, uncertain, or likely to cause disease. Specialists use laboratory evidence, family history, population data, and the person’s clinical findings.

\medterm{Varibaculum Cambriense} A bacterial species in the genus Varibaculum reported from human clinical specimens. Its role may be opportunistic and requires identification with the clinical context and susceptibility data.

\medterm{Variceal Band Ligation} A procedure using a flexible camera to place small bands around enlarged veins in the food pipe. The bands stop blood flow through the veins and are used to control or prevent bleeding.

\synonyms
Endoscopic Variceal Ligation; EVL

\medterm{Variceal Hemorrhage} Severe bleeding from enlarged, fragile veins in the food pipe or stomach, usually caused by high pressure in the liver’s blood vessels. It can cause vomiting blood, black stools, shock, and death without urgent treatment.

\medterm{Varicella} A primary infection with varicella-zoster virus that causes fever and a generalized itchy rash. See Chickenpox.

\synonyms
Chickenpox; Varicella-Zoster Virus Infection

\medterm{Varicella-Zoster Virus} The virus that causes chickenpox and shingles. After chickenpox, it can remain inactive in nerve tissue and reactivate years later as a painful, blistering rash called shingles.

\medterm{Varicellovirus} Varicellovirus is a genus of herpesviruses that includes varicella-zoster virus, the cause of chickenpox and shingles.

\medterm{Varices} Abnormally dilated and tortuous veins, commonly occurring in the legs or in the oesophagus and stomach as a result of portal hypertension.

\medterm{Varicocele} Enlargement of veins in the scrotum, often on the left. It may cause aching, a soft “bag of worms” feeling, testicular growth differences, or reduced fertility, but many cause no symptoms.

\medterm{Varicocele Embolization} A procedure that uses a thin tube to place coils or another material in enlarged scrotal veins. This blocks abnormal backward blood flow and reduces pressure in the veins.

\medterm{Varicose Ulcer} An open sore near the ankle caused by long-term poor drainage through the leg veins. It may be surrounded by swelling or skin discoloration and often needs compression and treatment of the underlying vein problem.

\medterm{Varicose Veins} Enlarged, twisted superficial veins, usually in the legs, caused by weak or damaged vein walls or valves that allow blood to pool and flow backward. They may cause visible bulging veins, aching, heaviness, itching, cramps, swelling, skin discoloration, bleeding, superficial inflammation or clotting, and difficult-to-heal venous ulcers.

\medterm{Varicose Veins and Venous Insufficiency} Enlarged superficial veins and inadequate return of blood from the legs because venous valves or the vein wall do not work effectively. Aching, heaviness, swelling, skin discoloration, dermatitis, and ulcers may occur.

\medterm{Variegate Porphyria} An inherited disorder of heme production that can cause attacks of abdominal pain, vomiting, nerve symptoms, confusion, or dark urine and can also cause fragile, blistering skin after sunlight exposure.

\medterm{Variola} A highly contagious and deadly disease caused by variola virus, eradicated globally by
1980 through vaccination. See Smallpox.

\medterm{Varix} An abnormally dilated, tortuous vein; multiple varices commonly occur in portal hypertension or chronic venous disease.

\medterm{Varus} A limb or joint deformity in which the part farther from the body bends inward, toward the body’s center line. Examples include bowlegs and some inward-turning foot deformities.

\medterm{Varus Stress Test} A clinical examination manoeuvre used to evaluate the lateral collateral ligaments of a joint, particularly the knee or elbow. A lateral-directed force is applied across the joint line while stabilizing the limb; increased laxity, gapping, or localized pain signifies injury to the lateral collateral ligament complex.

\medterm{Vas} A duct or vessel; in reproductive anatomy, the vas deferens transports sperm from the epididymis toward the ejaculatory duct.

\medterm{Vasa Deferentia} The paired muscular ducts that transport mature sperm from the epididymides toward the ejaculatory ducts during ejaculation.

\synonyms
Ductus Deferentes; Vas Deferens

\medterm{Vasa Efferentia} Small tubes that carry sperm from the testicle’s collecting channels to the epididymis, where sperm mature. They also remove some fluid and help move sperm onward.

\medterm{Vasa Previa} A pregnancy complication in which fetal blood vessels run across or near the cervix without the usual protection of the placenta or umbilical cord. If the membranes break, the vessels can tear and cause rapid fetal blood loss, so specialist planning is essential.

\medterm{Vasa Recta} Long, hairpin-shaped blood vessels in the inner part of the kidney. They run beside the loops of Henle and help the kidneys concentrate urine.

\medterm{Vasa Vasorum} Small blood vessels in the walls of large arteries and veins. They supply oxygen and nutrients to the outer parts of these vessel walls.

\medterm{Vascular} Relating to blood vessels or, sometimes, lymphatic vessels. Vascular conditions include narrowed or blocked arteries, aneurysms, blood clots, abnormal vessel connections, and problems with veins; the term may also describe tests, surgery, or medicines involving these vessels.

\medterm{Vascular Access} A route into a blood vessel for medicines, fluids, blood tests, monitoring, or dialysis. It may be a short-term catheter, a surgically joined artery and vein, or a synthetic graft, depending on the purpose and expected duration.

\medterm{Vascular Access For Hemodialysis} A route that allows blood to leave the body and return during dialysis. It may be a joined artery and vein, a synthetic graft, or a large vein catheter; each has different risks of infection, clotting, and failure.

\medterm{Vascular Balloon Catheter} An inflatable balloon mounted at the distal tip of a flexible vascular catheter used to dilate stenotic or occluded blood vessels during angioplasty, deliver and expand endovascular stents, or temporarily occlude vascular flow.

\synonyms
Angioplasty Balloon; Vascular Balloon Dilator; Balloon Catheter

\medterm{Vascular Bed} The network of blood vessels supplying a particular organ or tissue. Its vessels widen or narrow to adjust blood flow according to the tissue’s needs.

\medterm{Vascular Bypass} Surgery that uses a vein or artificial tube to carry blood around a blocked part of an artery and restore blood flow.

\synonyms
Surgical Bypass; Bypass Grafting; Arterial Bypass Graft

\medterm{Vascular Claudication} Alternate name; see \textit{Claudication}.

\medterm{Vascular Clip} A small device placed during a procedure to close or control a blood vessel. It may stop bleeding, prevent blood flow to a selected area, or secure a vessel permanently.
\synonyms
Hemostatic Clip; Surgical Vessel Clip; Ligation Clip

\medterm{Vascular Compliance} The ability of a blood vessel to stretch when pressure or blood volume increases. More compliant vessels can hold more blood with a smaller rise in pressure.

\synonyms
Angiopathic Complications; Vascular Adverse Events

\medterm{Vascular Complications} Problems involving blood vessels after disease, injury, or a procedure, such as bleeding, hematoma, thrombosis, embolism, vessel injury, or reduced blood flow to a limb. Symptoms and urgency depend on the vessel and tissue affected.

\medterm{Vascular Dementia} Problems with thinking, memory, mood, or walking caused by strokes or long-term damage to the brain’s small blood vessels. Symptoms may appear suddenly or gradually, and treatment focuses on preventing further vessel damage and supporting daily function.

\synonyms
VaD; Multi-Infarct Dementia; Vascular Cognitive Impairment

\medterm{Vascular Disease} Vascular disease affects arteries, veins, or lymphatic vessels and includes atherosclerosis, aneurysm, peripheral artery disease, venous thrombosis, varicose disease, and inflammatory disorders. Symptoms and treatment depend on the vessel and may involve risk-factor control, medicines, procedures, or surgery.

\medterm{Vascular Ehlers-Danlos Syndrome} A rare inherited connective-tissue disorder that makes arteries and hollow organs unusually fragile. It may cause thin translucent skin, easy bruising, characteristic facial features, or serious spontaneous blood-vessel or bowel rupture; specialist care and genetic counseling are important.

\synonyms
vEDS; Ehlers-Danlos Syndrome Type IV

\medterm{Vascular Embolic Particles} Small materials placed in a blood vessel during a procedure to block blood flow. They may be used to control bleeding or reduce blood flow to a tumor, fibroid, or other abnormal area.

\synonyms
Vascular Embolization Particle; PVA Particles

\medterm{Vascular Embolization} Intentional blockage of a blood vessel with material delivered through a catheter. It may stop bleeding, reduce blood flow to a tumor, or close an aneurysm or abnormal vessel connection.

\synonyms
Catheter Embolization; Therapeutic Embolization

\medterm{Vascular Embolization Coil} A catheter-delivered wire device used to block selected blood vessels, often to treat aneurysms or bleeding.

\synonyms
Endovascular Coil

\medterm{Vascular Embolization Plug} A vascular occlusion device used to block a selected vessel under image guidance. A device placed in a blood vessel to block selected blood flow during image-guided treatment.
\synonyms
Vascular Occlusion Plug; Embolization Plug; Amplatzer Plug

\medterm{Vascular Endothelial Growth Factor} A protein that helps new blood vessels form and makes blood vessels more leaky. Some medicines block it to treat certain eye diseases and cancers.
\synonyms
VEGF; Vascular Permeability Factor

\medterm{Vascular Endothelial Growth Factor Inhibitor} A medicine that blocks vascular endothelial growth factor, a protein that helps blood vessels grow. These medicines are used for some cancers and eye diseases involving abnormal new blood vessels.

\synonyms
Anti-VEGF Agent

\medterm{Vascular Endothelial Growth Factor Inhibitors} Medicines that block vascular endothelial growth factor and reduce abnormal blood-vessel growth or leakage. They are used for some eye diseases and cancers.

\synonyms
Anti-VEGF Agents; VEGF Neutralizers

\medterm{Vascular Forceps} Surgical forceps designed to handle blood vessels, clots, tissue, or medical materials during vascular procedures. Appropriate use helps avoid crushing or tearing delicate vessels.

\synonyms
Vascular Surgical Forceps; Atraumatic Clamp

\medterm{Vascular Graft} A tube used to replace or bypass a damaged or blocked blood vessel. It may be made from a person's vein or from a synthetic material; possible problems include infection, bleeding, narrowing, or a clot blocking the graft.

\medterm{Vascular Graft Patency} The continued openness and useful blood flow through a blood-vessel graft or bypass. A clot, narrowing, infection, or scar growth can block it, so examination and sometimes ultrasound are used for monitoring.

\medterm{Vascular Hemostat} A surgical clamp used to temporarily occlude a blood vessel and control bleeding.
\synonyms
Vascular Clamp; Atraumatic Hemostasis Forceps

\medterm{Vascular Imaging} Tests that show blood vessels and sometimes measure blood flow. Ultrasound, CT, MRI, or catheter-based imaging can find narrowing, blockage, aneurysm, dissection, clot, bleeding, or abnormal vessels; the safest test depends on the body area, urgency, kidney function, and contrast risks.

\medterm{Vascularity} The amount and arrangement of blood vessels supplying a tissue or growth. A highly vascular area may heal or bleed differently and may look different on imaging; blood supply can also affect some treatments.

\medterm{Vascular Laser Therapy} A treatment that uses laser energy to close or change selected blood vessels. The laser may be applied through the skin or through a thin tube.
\synonyms
Pulsed Dye Laser; PDL; Vascular Laser Therapy

\medterm{Vascular Ligation} Surgical tying, clipping, or sealing of a blood vessel to stop bleeding or redirect blood flow. It may be used during surgery or to treat selected abnormal vessels.
\synonyms
Vascular Occlusion; Vessel Clipping

\medterm{Vascular Malformation} A blood-vessel or lymph-vessel structure that developed abnormally before birth. It may cause a skin mark, swelling, pain, bleeding, or pressure on nearby organs; effects depend on its type and location.

\medterm{Vascular Needle Holder} A surgical instrument used to hold and guide fine suture needles during repair of blood vessels. It permits precise stitches while minimizing crushing or tearing of delicate vessel walls.
\synonyms
Vascular Needle Holder; Castroviejo Needle Holder

\medterm{Vascular Pain} Pain caused by impaired blood flow, often appearing in a limb with exertion as claudication or, in more severe disease, at rest.

\medterm{Vascular Patch Graft} A piece of biological or synthetic material sewn onto a blood vessel to widen or repair it.

\synonyms
Vascular Patch Graft; Angioplasty Patch

\medterm{Vascular Retraction} To pull a vessel or tissue away from an operative field using a surgical retractor.
\synonyms
Vessel Retraction; Surgical Vessel Exposure

\medterm{Vascular Ring} A birth defect in which blood vessels around the heart form a ring that can press on the windpipe or food pipe. It may cause noisy breathing, cough, repeated chest infections, swallowing difficulty, or feeding problems. Some people need no treatment; symptomatic rings may require surgery.

\medterm{Vascular Scissors} A surgical instrument with two blades used to cut tissue, sutures, dressings, or other materials
during vascular procedures.

\synonyms
Vascular Scissors; Potts Scissors

\medterm{Vascular Shunt} A surgically created or implanted pathway that redirects blood or another fluid from one vessel or body space to another.

\synonyms
Vascular Shunt; Bypass Shunt

\medterm{Vascular Sponge} A sterile absorbent material used during blood-vessel procedures to absorb blood, apply pressure, or protect delicate tissue.
\synonyms
Vascular Laparotomy Sponge; Surgical Cottonoid

\medterm{Vascular Staple} A small metal or polymer device used to close tissue or join selected blood-vessel structures. It provides rapid closure but can cause narrowing, bleeding, thrombosis, or tissue injury if unsuitable or misplaced.
\synonyms
Vascular Surgical Staple; Vessel Closure Device

\medterm{Vascular Stent} A small mesh or tube placed inside a narrowed or blocked vessel to help keep it open.

\synonyms
Endovascular Stent; Vascular Prosthesis

\medterm{Vascular Stent Balloon} A balloon on a catheter used to expand a vascular stent against the vessel wall and restore an open channel. Inflation can cause vessel injury, bleeding, clotting, or rarely rupture.
\synonyms
Stent-Expansion Balloon; Angioplasty Balloon

\medterm{Vascular Stent-Graft} A fabric-covered metal frame placed inside a blood vessel to reinforce a weak wall or seal an aneurysm, tear, or leak. Follow-up imaging checks for movement, blockage, or leakage around the graft.
\synonyms
Endovascular Graft; Aortic Stent-Graft

\medterm{Vascular Suction} A suction device used during a vascular procedure to remove blood or other fluid from the operative
field.

\synonyms
Vascular Suction Cannula; Frazier Suction

\medterm{Vascular Surgery} Operations that repair, widen, replace, or reroute blood vessels to restore circulation or prevent bleeding. Procedures may treat aneurysms, blocked arteries, blood-vessel injuries, or problems with veins and dialysis access.

\medterm{Vascular System} The network of arteries, veins, and capillaries that carries blood throughout the body. It delivers oxygen, nutrients, and hormones to tissues and carries carbon dioxide and other waste products away. Disease may narrow, weaken, block, or inflame these vessels.

\medterm{Vascular Thrombectomy} Removal of a blood clot from a blood vessel using surgery or a catheter-based procedure. The aim is to restore blood flow and protect tissue from damage when a clot blocks an important artery or vein.

\synonyms
Vascular Thrombectomy; Clot Extraction

\medterm{Vascular Tissue Adhesive} A liquid material used to block selected blood vessels during image-guided embolization. Doctors inject it through a catheter to control bleeding, close an abnormal vessel, or reduce blood flow to a tumor.

\synonyms
Cyanoacrylate Embolization Glue; Endovascular Embolic Adhesive

\medterm{Vascular Tone} The usual degree of tightening in the muscle within blood-vessel walls. It helps control vessel width, blood flow, and blood pressure.

\medterm{Vascular Trauma} Injury to an artery or vein caused by a blunt blow, penetrating wound, crush injury, or medical procedure. It may cause severe bleeding or loss of blood flow and needs urgent assessment of circulation and nearby nerves.

\medterm{Vascular Ulcer} A persistent open wound caused or maintained by poor blood flow, high pressure in the veins, loss of protective feeling, or another blood-vessel problem. It usually affects the legs or feet and may heal slowly or become infected.

\synonyms
Vascular Ulceration; Ischemic Ulcer; Arterial Ulcer

\medterm{Vasculitic Neuropathy} Nerve damage caused by inflammation of the small blood vessels that supply nerves. It often causes painful, uneven numbness or weakness and may affect several limbs as the vessel inflammation continues.

\medterm{Vasculitis} Inflammation of blood vessels that can narrow them, reduce blood flow, or cause bleeding and tissue damage. It may affect large, medium, or small vessels and can involve the skin, nerves, kidneys, lungs, or other organs.

\synonyms
Angiitis; Arteritis; Blood Vessel Inflammation

\medterm{Vas Deferens} A muscular tube that carries sperm from the epididymis toward the urethra during ejaculation.

\medterm{Vasectomy} A permanent contraceptive procedure in which the vas deferens are divided or occluded to prevent sperm from entering the ejaculate. It does not immediately provide sterility and does not protect against sexually transmitted infections.

\medterm{Vasectomy Reversal} Surgery to reconnect the tubes that carry sperm after a vasectomy. It may restore sperm to the semen, but pregnancy is not guaranteed.
\synonyms
Vasovasostomy; Vasoepididymostomy

\medterm{Vasoactive Intestinal Peptide Test} A blood test measuring a hormone that affects the intestine and blood vessels. It is used mainly when a rare hormone-producing tumor is suspected in someone with severe watery diarrhea, dehydration, or abnormal body salts.

\medterm{Vasoconstriction} Narrowing of blood vessels caused by contraction of vascular smooth muscle, which can raise blood pressure or reduce blood flow to affected tissues.

\medterm{Vasoconstrictor} A substance that narrows blood vessels by making their muscle walls contract. This can raise blood pressure and reduce local blood flow; vasoconstrictor medicines are used in selected emergencies or with local anesthetics and require careful monitoring.

\medterm{Vasoconstrictors} Substances or medicines that narrow blood vessels by making their muscular walls contract. They can raise blood pressure or reduce local blood flow and are used only for selected conditions because excessive narrowing can be harmful.

\medterm{Vasodepressor Syncope} Fainting caused mainly by sudden widening of blood vessels and a fall in blood pressure. It often occurs as part of a vasovagal reaction and may be preceded by nausea, sweating, warmth, or lightheadedness.

\medterm{Vasodilatation} Widening of blood vessels caused by relaxation of vascular smooth muscle, increasing blood flow and potentially lowering blood pressure.

\medterm{Vasodilation} Widening of blood vessels when their muscular walls relax. It can increase blood flow and lower blood pressure and occurs with exercise, heat, medicines, and inflammation.

\medterm{Vasodilator} A substance that relaxes blood-vessel walls and makes the vessels wider. It can lower blood pressure or improve blood flow and is used for selected heart, blood-pressure, or circulation problems; excessive widening may cause headache, flushing, dizziness, or fainting.

\medterm{Vasodilators} Medicines or substances that relax vascular smooth muscle and widen blood vessels. They are used for conditions such as hypertension, angina, heart failure, or pulmonary hypertension depending on the agent.

\medterm{Vasodilatory Shock} A form of distributive shock characterized by severe systemic arterial hypotension and tissue hypoperfusion resulting from profound, refractory peripheral vasodilation and loss of vascular tone despite preserved or elevated cardiac output, most frequently seen in severe sepsis, anaphylaxis, and neurogenic injury.

\medterm{Vasomotor} Relating to control of blood-vessel width by nerves, hormones, and local chemicals. Vasomotor changes narrow or widen vessels and help regulate blood pressure, body temperature, tissue blood flow, and skin color.

\medterm{Vasomotor Center} A group of nerve cells in the lower brain that helps control the widening and narrowing of blood vessels and helps maintain blood pressure.

\medterm{Vasomotor Centre} Groups of nerve cells in the brainstem that help control blood-vessel narrowing and widening. By changing vessel tone, they influence blood pressure, heart and tissue blood flow, and the body’s response to standing or stress.

\medterm{Vasomotor Nerves} Autonomic nerves controlling blood-vessel widening and narrowing by changing smooth-muscle activity, helping regulate blood pressure and temperature.

\medterm{Vasomotor Rhinitis} Common subtype term; see \textit{Nonallergic rhinitis}.

\medterm{Vaso-Occlusive Crisis} The classic acute, excruciatingly painful hallmark complication of sickle cell disease, caused by the intravascular sludging and adherence of rigid, polymer-filled sickled erythrocytes in the microvasculature, producing focal tissue ischemia, infarction, and systemic inflammation in bones, joints, or viscera.

\synonyms
Sickle Cell Pain Crisis; VOC

\medterm{Vasopressin} A hormone that increases water reabsorption in the kidney and can constrict
blood vessels. It is released from the posterior pituitary.

\synonyms
Arginine Vasopressin; AVP

\medterm{Vasopressin Receptor} A cell sensor that responds to the hormone vasopressin. Different types help the kidneys retain water and help blood vessels narrow, influencing body-fluid balance and blood pressure.

\medterm{Vasopressin Receptor Antagonists} A pharmacological class of aquaretic medications (vaptans, such as tolvaptan and conivaptan) that competitively block vasopressin V2 receptors in the renal collecting duct, promoting solute-free water excretion without electrolyte loss in euvolemic and hypervolemic hyponatremia (SIADH, heart failure, cirrhosis).

\synonyms
Vaptans

\medterm{Vasopressor} A medicine or body signal that raises blood pressure, mainly by narrowing blood vessels and sometimes by strengthening the heartbeat. These medicines are used in selected emergencies such as shock and require close monitoring.

\medterm{Vasopressors In Anesthesia} Medicines that increase vascular tone or cardiac output to support blood pressure during
anesthesia or critical illness. The choice and dose depend on the cause of hypotension,
heart rate, cardiac function, and local protocol; monitoring is required.

\medterm{Vasospasm} Sudden narrowing of a blood vessel caused by contraction of its muscular wall. It can temporarily reduce blood flow and cause chest pain, headache, stroke-like symptoms, or tissue injury, depending on the vessel affected.

\medterm{Vasovagal Attack} A transient reflex episode caused by reduced sympathetic tone and increased vagal activity, producing hypotension and often bradycardia with fainting or near-fainting. Common triggers include prolonged standing, pain, fear, heat, dehydration, or emotional stress.

\medterm{Vasovagal Reaction} A reflex that slows the heart and widens blood vessels, briefly lowering blood pressure and sometimes causing fainting; triggers include pain, fear, prolonged standing, heat, or blood exposure.

\medterm{Vasovagal Syncope} A common faint caused by a reflex that briefly slows the heart or widens blood vessels, lowering blood pressure and brain blood flow. Triggers include standing still, heat, pain, fear, or seeing blood; warning symptoms often occur first.

\synonyms
Neurocardiogenic Syncope

\medterm{Vater-Pacini Corpuscles} Large, rapidly adapting, encapsulated mechanoreceptors located deep in the dermis, subcutaneous tissue, periosteum, and joint capsules. Morphologically composed of a central unmyelinated nerve axon surrounded by multiple concentric lamellae of flattened Schwann-like cells separated by fluid; specialized physiologically to detect high-frequency mechanical vibration (200 to 300 Hz) and transient acceleration.

\synonyms
Pacinian Corpuscles; Lamellar Corpuscles

\medterm{VATS} Alternate name; see \textit{Video-Assisted Thoracoscopic Surgery}.

\medterm{Vaughan Williams Classification} A system that groups medicines used for abnormal heart rhythms according to their main effect on the heart’s electrical activity. It helps describe treatment but does not replace patient-specific prescribing.

\medterm{VCJD} Alternate name; see \textit{Creutzfeldt-Jakob disease}.

\medterm{VDRL Test} A blood or spinal-fluid test that looks for an immune reaction associated with syphilis. A reactive result needs a more specific confirmatory test and may help monitor response to treatment.

\medterm{Vector} A living carrier, usually an insect or tick, that spreads an infection between people or animals. Some vectors allow the germ to develop inside them; others transfer it mechanically without the germ changing.

\medterm{Vector-Borne Disease} An infection spread by a living carrier such as a mosquito, tick, flea, or sandfly. Examples include malaria, dengue, Lyme disease, and some forms of viral encephalitis; prevention focuses on avoiding bites and controlling the carrier.

\medterm{Vector Control} Measures that reduce insects or other animals that spread infection. Examples include removing standing water, using insecticides, treating breeding areas, improving buildings, wearing protective clothing, and using repellents or bed nets.

\medterm{Vegetarian} Vegetarian describes a diet that avoids meat and may include dairy products, eggs, or both. Nutritional needs can usually be met with varied foods, but iron, vitamin B12, protein, and other nutrients may need attention.

\medterm{Vegetation} A mass that can form on a heart valve during infective endocarditis. It may contain bacteria, blood-clotting material, and immune cells and can damage the valve or break off and travel to other organs.

\medterm{Vein} A blood vessel that carries blood toward the heart. Most veins carry oxygen-poor blood, but veins from the lungs and placenta carry oxygen-rich blood. Veins often contain valves that help prevent backward flow.

\medterm{Vein-Graft Aneurysms} An abnormal widening of a vein used to bypass a blocked heart artery. It may contain a clot, block blood flow, press on nearby structures, or rupture and therefore needs specialist assessment.

\medterm{Vein of Galen Malformation} A rare congenital intracranial vascular anomaly resulting from an embryonic arteriovenous fistula that drains directly into the primitive median prosencephalic vein (the precursor of the great cerebral vein of Galen), presenting in neonates with high-output congestive heart failure and hydrocephalus.

\synonyms
VGM; Vein of Galen Aneurysmal Malformation

\medterm{VeIP Regimen} A cancer-treatment combination of vinblastine, ifosfamide, and cisplatin, used mainly for selected germ-cell tumors. It can lower blood counts and damage the kidneys, nerves, or bladder, so specialist monitoring is required.

\medterm{Velamentous Cord Insertion} An abnormal attachment in which the umbilical blood vessels travel through the membranes before reaching the placenta without the usual protective tissue. The vessels can be compressed or torn, increasing risks of poor fetal growth, bleeding, or vasa previa and requiring specialist pregnancy monitoring.

\synonyms
Velamentous Insertion of the Umbilical Cord

\medterm{Velocardiofacial Syndrome} Alternate name; see \textit{DiGeorge syndrome (22q11.2 deletion syndrome)}.

\medterm{Velopharyngeal Insufficiency} A structural or functional anatomical failure of the soft palate (velum) and posterior pharyngeal wall to achieve complete closure during speech and deglutition, resulting in hypernasal speech resonance, nasal air emission, and nasal regurgitation of fluids.
\synonyms
VPI; Velopharyngeal Incompetence

\medterm{Vemurafenib} A targeted cancer medicine for tumors with a specific BRAF V600 gene change, including some melanomas. It can cause joint pain, sun sensitivity, abnormal heart rhythm, new skin cancers, or treatment resistance.

\medterm{Vena Cava} Either of the two largest veins returning blood to the right side of the heart. The superior vena cava drains the upper body, and the inferior vena cava drains the abdomen, pelvis, and legs.


\synonyms
Venae Cavae; Superior and Inferior Vena Cava

\medterm{Vena Cava Filter} A small device placed in the large vein leading to the heart to trap blood clots traveling from the legs toward the lungs. It may be used when blood-thinning medicines cannot be used or have not worked.

\synonyms
IVC Filter; Greenfield Filter

\medterm{Venepuncture} Inserting a needle into a vein to collect blood or place intravenous access. It may cause brief pain, bruising, bleeding, infection, or fainting, and clean technique reduces complications.

\medterm{Venereal Disease Research Laboratory Test} A blood test that looks for antibodies associated with syphilis. It is used for screening and monitoring treatment, but a reactive result needs confirmation with a treponemal test.

\medterm{Venesection} The controlled removal of blood from a vein, used diagnostically or therapeutically, for example in selected cases of hemochromatosis or polycythemia.

\medterm{Venipuncture} Puncture of a vein with a needle to obtain blood, administer fluids or medicines, or establish intravenous access. Bruising, infection, bleeding, and nerve or artery injury are uncommon complications.

\medterm{Venogram} Contrast imaging of veins used to detect thrombosis, obstruction, abnormal anatomy, or selected venous flow disorders.

\medterm{Venography} Imaging of veins, usually after contrast material is administered, to assess venous anatomy, obstruction, thrombosis, or malformation.

\medterm{Venom} A poisonous substance delivered by an animal through a bite, sting, or similar structure. Its effects vary by species and may include pain, swelling, nerve problems, bleeding, or dangerous changes in breathing or blood pressure.

\medterm{Venomous Bite} A bite that injects venom from an animal such as a snake, spider, or arthropod. Effects range from local
pain and swelling to bleeding, paralysis, or systemic toxicity.

\medterm{Veno-Occlusive Disease} Blockage and injury of small veins in the liver, especially after some chemotherapy or stem-cell transplantation. It can cause painful liver enlargement, rapid weight gain, fluid retention, and jaundice and can be life-threatening.

\medterm{Veno-occlusive Hepatic Disease (Sinusoidal Obstruction Syndrome)} A potentially severe liver injury caused by narrowing or blockage of hepatic sinusoids, classically after hematopoietic-cell transplantation or certain chemotherapy. Painful liver enlargement, weight gain, fluid retention, jaundice, and rising bilirubin may occur.

\medterm{Venous} Relating to veins, which usually carry blood back toward the heart. Vein problems include blood clots, varicose veins, poor return of blood, and ulcers; pulmonary veins are an exception because they carry oxygen-rich blood from the lungs to the heart.

\medterm{Venous Air Embolism} Obstruction of the venous circulation by intravascular gas, usually after procedures, trauma, or pressure changes. It can impair pulmonary or cardiac function and, if gas reaches the arterial circulation, cause neurologic or systemic ischemia.

\medterm{Venous Aneurysm} A localized abnormal dilatation of a vein. It may be asymptomatic or present as a soft mass, pain, thrombosis, embolism, or bleeding depending on its location and size.

\medterm{Venous Blood Gas} The clinical laboratory measurement of pH, partial pressure of carbon dioxide ($p\text{CO}_2$), bicarbonate ($\text{HCO}_3^-$), and base excess in a peripheral or central venous blood sample to assess systemic acid-base balance, ventilatory status, and tissue perfusion.

\synonyms
Venous Blood Gas; VBG; Central Venous Blood Gas
\medterm{Venous Capacitance} The physical and physiological capacity of the systemic venous circulation to accommodate large changes in blood volume with minimal alterations in venous pressure, functioning as the primary high-capacitance, low-resistance reservoir holding approximately 60 to 70 percent of total blood volume at rest.

\medterm{Venous Diseases} Conditions affecting veins, including blood clots, varicose veins, poor return of blood, inflammation, ulcers, injuries, and birth-related malformations. They may cause pain, swelling, skin changes, bleeding, or dangerous clots traveling to the lungs.

\medterm{Venous Hum} A harmless continuous humming sound heard over the neck veins, most often in healthy children. It changes or disappears with body position or gentle pressure on the vein, helping distinguish it from heart murmurs that need further evaluation.

\synonyms
Cervical Venous Hum; Innocent Venous Hum

\medterm{Venous Insufficiency} Impaired return of blood through the veins, usually in the legs, because valves allow backward flow, veins are blocked, or the calf-muscle pump is ineffective. Blood may pool and cause aching, heaviness, swelling, varicose veins, itching, darkened or hardened skin, stasis dermatitis, and difficult-to-heal ulcers near the ankle.

\medterm{Venous Leak} Difficulty maintaining an erection because blood leaves the penis too readily during erection. It may reflect vascular, nerve, hormonal, or medication-related problems and should be evaluated as erectile dysfunction.
\synonyms
Veno-Occlusive Dysfunction; Cavernous Venous Leakage

\medterm{Venous Pressure} The pressure of blood inside the veins. It depends on blood volume, vein width, gravity, and how well the right side of the heart pumps; abnormally high pressure can cause swelling, enlarged veins, or fluid buildup.

\medterm{Venous Return} The flow of blood through the veins back to the heart. It is helped by vein valves, muscle movement, breathing, and pressure changes in the chest.

\medterm{Venous Stasis Dermatitis} Inflammation of the skin of the lower legs caused by chronic swelling and pooling of blood and fluid, usually from chronic venous insufficiency. It may cause itching, redness, scaling, skin darkening, thickening, or ulcers; treatment addresses the swelling and protects the skin.

\medterm{Venous Stasis Ulcer} A long-lasting wound near the ankle caused by poor blood return through the leg veins. The leg may also be swollen, painful, or have darkened, thickened skin; compression and treatment of the vein problem may help.

\medterm{Venous Thromboembolism} A blood clot that forms in a vein, usually deep in the leg, and may travel to the lungs. A clot in the lungs can cause sudden breathlessness, chest pain, collapse, or death.

\medterm{Venous Thromboembolism In Pregnancy} A blood clot in a deep vein or lungs during pregnancy or after birth. Pregnancy increases clot risk, and symptoms such as one-sided leg swelling, chest pain, sudden breathlessness, or coughing blood can indicate a serious clot.

\medterm{Venous Thromboembolism Prophylaxis} Measures used to prevent dangerous vein clots, especially around surgery or hospitalization. They may include walking early, compression devices or stockings, and blood-thinning medicine when bleeding risk is acceptable.

\medterm{Venous Thrombosis} The formation of a blood clot (thrombus) within a vein, most commonly in the deep veins
of the lower extremities (deep vein thrombosis, DVT).

\medterm{Venous Ulcer} A chronic wound caused by impaired venous drainage, usually occurring near the ankle in a limb with chronic venous insufficiency.

\medterm{Ventilation} Movement of air into and out of the lungs; effective ventilation brings oxygen to the alveoli and removes carbon dioxide.

\medterm{Ventilation-Perfusion Mismatch} A mismatch between air reaching parts of the lungs and blood reaching those same parts. It can lower blood oxygen and occur with pneumonia, COPD, asthma, heart failure, or a lung clot.

\medterm{Ventilation-Perfusion Ratio} The relationship between air reaching the lungs and blood reaching the lung tissues. Good matching allows efficient oxygen exchange; poor matching can cause low blood oxygen.

\medterm{Ventilation-Perfusion Scan} A nuclear-medicine scan that compares air movement in the lungs with blood flow through them. Areas with blood flow blocked but preserved air movement may suggest a pulmonary embolism. The test uses a breathed tracer and an injected tracer and is interpreted with symptoms and other findings.

\medterm{Ventilator} A machine that supports or replaces spontaneous breathing by delivering air, oxygen, or an air-oxygen mixture under positive pressure. Noninvasive ventilators deliver support through a mask, while invasive ventilators connect through an endotracheal or tracheostomy tube. Ventilators can provide different patterns and levels of pressure or volume according to the device and clinical purpose.

\synonyms
Mechanical Ventilator
Breathing Machine

\medterm{Ventilator-Associated Pneumonia} A lung infection that develops after a person has used a breathing machine for at least two days. It may cause fever, worsening oxygen levels, cough, or abnormal lung findings and requires prompt evaluation and treatment.

\medterm{Ventilator-Induced Lung Injury} Lung damage caused or worsened by a breathing machine when pressure, volume, or oxygen levels are excessive. Clinicians use protective settings and close monitoring to support breathing while limiting further injury.

\medterm{Ventouse} A vacuum extractor placed on the baby's scalp to assist a vaginal birth when delivery needs help. It is used only when specific safety conditions are met and can cause temporary scalp swelling or bruising.

\medterm{Ventral} Relating to the belly or front side of the body; synonym for anterior in human anatomy.

\medterm{Ventral Body Cavity} The body cavity toward the front. It includes the thoracic cavity and the abdominopelvic cavity, which contain many of the body’s major organs.

\medterm{Ventral Hernia} Protrusion of abdominal tissue through a weakness in the front abdominal wall, including incisional and umbilical hernias. It may cause a bulge or pain and requires urgent assessment if incarceration or strangulation is suspected.

\medterm{Ventral Root} The front root of a spinal nerve. It carries movement and other outgoing signals from the spinal cord.

\medterm{Ventricle} One of the two lower chambers of the heart that pump blood out of it. The right ventricle sends blood to the lungs to receive oxygen; the left ventricle sends oxygen-rich blood to the rest of the body.

\medterm{Ventricular} Relating to a ventricle. In the heart, it concerns one of the lower pumping chambers; in the brain, it concerns the fluid-filled spaces that contain cerebrospinal fluid. The organ gives the term its specific meaning.

\medterm{Ventricular Aneurysm} A localized, thin, fibrous, dyskinetic or akinetic bulging outpouching of the scarred left ventricular wall developing as a late mechanical complication of transmural myocardial infarction (typically anterior MI), predisposing to heart failure, refractory ventricular arrhythmias, and apical mural thrombus formation.

\medterm{Ventricular Arrhythmia} An abnormal heart rhythm arising from the ventricles. It may cause palpitations, dizziness, fainting,
chest discomfort, or sudden cardiac arrest; fainting or severe symptoms require emergency care.

\medterm{Ventricular Arrhythmias} Abnormal heart rhythms arising from the ventricles. They range from premature ventricular contractions to ventricular tachycardia and fibrillation, which can cause sudden cardiac arrest.

\medterm{Ventricular Assist Device} A mechanical pump that helps a weakened heart move blood. It may support the left side, right side, or both and can be temporary or long-term, but bleeding, infection, stroke, and clots are important risks.

\medterm{Ventricular Bigeminy} A cardiac rhythm pattern on electrocardiogram in which every normal sinus beat is consistently followed by a premature ventricular complex (PVC) in a regular alternating 1:1 sequence, commonly provoked by electrolyte derangements, digitalis toxicity, or underlying ischemic heart disease.

\medterm{Ventricular Fibrillation} A chaotic, disorganized ventricular rhythm in which the ventricles quiver rather than contract, producing no effective cardiac output and cardiac arrest. It is a shockable rhythm; immediate CPR and unsynchronized defibrillation are required, with treatment of the underlying cause.

\synonyms
V-Fib

\medterm{Ventricular Hypertrophy} Thickening of the wall of one of the heart’s lower chambers. It may develop from long-term high blood pressure, valve disease, inherited heart muscle disease, or intense training and can lead to stiffness, abnormal rhythms, or heart failure.

\medterm{Ventricular Myocardium} The heart muscle forming the walls of the ventricles. Its contraction pumps blood to the lungs and the rest of the body.

\medterm{Ventricular Non-Compaction} A primary genetic cardiomyopathy characterized by a thickened myocardial wall with a non-compacted endocardial layer displaying prominent, deep intertrabecular recesses and deep trabeculations, typically affecting the left ventricular apex and lateral wall.

\synonyms
LVNC; Non-Compaction Cardiomyopathy

\medterm{Ventricular Premature Contraction} Alternate name; see \textit{Premature ventricular contractions (PVCs)}.

\medterm{Ventricular Remodeling} Progressive alterations in left ventricular geometry, mass, chamber volume, and spherical shape occurring after acute myocardial infarction or chronic pressure/volume overload, driven by neurohormonal activation, cardiomyocyte hypertrophy, and interstitial fibrosis.
\synonyms
Post-Infarction Remodeling; Cardiac Remodeling

\medterm{Ventricular Rupture} A tear in a wall of one of the heart’s lower chambers, often after a heart attack. It can cause life-threatening bleeding around the heart.

\medterm{Ventricular Septal Defect} A hole in the wall between the heart’s lower chambers. Blood may flow from the stronger left side to the right and increase blood flow to the lungs. Small holes may cause no problems; larger ones can cause heart failure or lung-vessel damage.

\medterm{Ventricular Septal Defect Cross-Reference} Alternate name; see \textit{Ventricular septal defect}.

\medterm{Ventricular Tachycardia} A rapid heart rhythm that starts in the lower chambers of the heart. It may cause palpitations, dizziness, fainting, shock, or cardiac arrest and can be life-threatening.

\medterm{Ventriculoatrial} Relating to a connection between a ventricle and an atrium. In brain surgery, a ventriculoatrial shunt is a tube that drains extra fluid from a brain cavity into the right upper chamber of the heart.

\medterm{Ventriculography} Imaging used to measure how well the heart’s ventricles pump blood. A radionuclide scan or cardiac MRI can estimate the amount of blood ejected with each heartbeat and show abnormal wall movement; the method is chosen for the clinical question.

\medterm{Ventriculoperitoneal Shunt} A surgically implanted system that drains cerebrospinal fluid from a brain ventricle into the abdominal cavity, where the fluid
can be absorbed.

\medterm{Ventrosuspension} Surgical fixation of a displaced or prolapsed uterus to the anterior abdominal wall or nearby supporting tissues.

\medterm{Venturi Mask} A high-flow oxygen delivery device that uses the Venturi principle to supply a precise, fixed fraction of inspired oxygen ($\text{FiO}_2$, typically 24\% to 50\%). Oxygen passing through a calibrated narrow jet nozzle creates a pressure drop that draws in a predictable volume of room air through color-coded entrainment ports. Because the total gas flow meets or exceeds the patient's peak inspiratory flow rate, the delivered oxygen concentration remains constant regardless of the patient's breathing rate or depth, making it vital for managing hypoxemia in patients with chronic obstructive pulmonary disease at risk of hypercapnic respiratory failure.
\synonyms{Air-Entrainment Mask; Venturi Oxygen Mask; High-Flow Oxygen Mask}

\medterm{Venturi Principle} A physical principle in fluid dynamics stating that as a fluid or gas flows through a constricted section of a pipe or nozzle, its velocity increases and its static lateral pressure decreases. In medicine, this localized pressure drop creates suction that entrains ambient air or fluids, forming the operational basis for air-entrainment oxygen masks (Venturi masks), jet nebulizers, suction aspirators, and explaining hemodynamic phenomena such as systolic anterior motion of the mitral valve in hypertrophic obstructive cardiomyopathy.
\synonyms{Venturi Effect; Bernoulli-Venturi Principle; Air-Entrainment Principle}

\medterm{Venule} A small blood vessel that collects blood from capillaries and carries it into larger veins. Venules also allow white blood cells and fluid to leave the bloodstream during inflammation.

\medterm{Venules} Small blood vessels that collect blood from capillaries and join to form veins. Their walls also participate in inflammation by allowing white blood cells to leave the bloodstream.

\medterm{Verapamil} A non-dihydropyridine calcium-channel blocker used for selected hypertension, angina, and supraventricular tachyarrhythmias; it can cause bradycardia and constipation.

\medterm{Verapamil Toxicity} Acute severe poisoning from the non-dihydropyridine calcium channel blocker verapamil, characterized by profound peripheral vasodilation, profound myocardial depression, high-grade atrioventricular block, refractory bradycardia, and cardiogenic shock; managed with intravenous calcium, high-dose insulin euglycemia therapy, and lipid emulsion.

\medterm{Vermicides} Medicines or chemicals that kill parasitic worms. The correct treatment depends on the worm, the body site affected, the person’s age, and other health conditions.

\medterm{Vernal Keratoconjunctivitis} A severe, recurrent allergic eye inflammation that primarily affects young boys in warm, sunny climates and often flares up during the spring and summer (spring catarrh). It causes intense eye itching, light sensitivity, watering, and thick, stringy mucus. Characteristic findings include large, cobblestone-like bumps on the underside of the upper eyelid (giant papillae) and small white dots along the border of the cornea (Horner-Trantas dots). In severe cases, inflammation and eye-rubbing can create a non-healing sore on the clear front of the eye (a shield ulcer), which can lead to corneal scarring and vision changes. The condition typically improves on its own after puberty.

\synonyms
VKC; Spring Catarrh; Vernal Catarrh; Vernal Conjunctivitis

\medterm{Verner-Morrison Syndrome} A rare neuroendocrine tumor syndrome (WDHA syndrome) caused by a vasoactive intestinal peptide-secreting pancreatic islet cell tumor (VIPoma), characterized by the tetrad of severe watery diarrhea ('pancreatic cholera'), hypokalemia, achlorhydria (hypochlorhydria), and metabolic acidosis.
\synonyms
VIPoma Syndrome; WDHA Syndrome

\medterm{Vernix} A white, greasy protective coating on the skin of a fetus or newborn, composed mainly of water, lipids, and shed skin cells.

\medterm{Vernix Caseosa} A white, greasy protective coating on the skin of a fetus or newborn, composed mainly of water, lipids, and shed skin cells.

\medterm{Verruca} Alternate name; see \textit{Wart}.

\medterm{Verruca Plana} Small, smooth, relatively flat warts caused by HPV, often occurring in groups on the face, hands, or legs.

\synonyms
Flat Warts; Juvenile Warts

\medterm{Verrucas} Verrucas are rough skin growths on the feet caused by human papillomavirus. They may be painful with pressure and often resolve without treatment, although treatment can reduce discomfort or spread.

\medterm{Verruca Vulgaris} A common, benign, hyperkeratotic, exophytic cutaneous papilloma with a rough, verrucous ('warty') surface and thrombosed capillary black dots, caused by localized epidermal infection by human papillomavirus (HPV types 2, 4, 7, and 1).
\synonyms
Common Wart; Viral Verruca

\medterm{Verruciform Xanthoma} A rare, noncancerous growth usually found in the mouth or genital area. It may look like a wart, so examination and sometimes a tissue sample are needed to distinguish it from cancer.

\medterm{Verrucose} Having a rough, thickened, wart-like, or finger-like surface. It is a descriptive term for a skin, mucosal, or other lesion and does not by itself identify whether the lesion is benign, infectious, or cancerous.

\medterm{Verrucous} Having a rough, thickened, irregular, or wart-like surface. The appearance can occur in benign growths, infections, chronic inflammation, or some cancers.

\medterm{Verrucous Carcinoma} A low-grade, well-differentiated variant of squamous cell carcinoma with a warty,
exophytic appearance. It is locally invasive but rarely metastasizes, commonly affecting
the oral cavity, skin, or genitalia.

\medterm{Version} A change in the position or presentation of a body structure. In pregnancy, fetal version is a procedure or movement that changes the baby’s position, often to place the head downward before birth.

\medterm{Vertebra} One of the bones forming the spinal column. Each vertebra helps support the body, allows movement, provides attachment for muscles, and surrounds and protects the spinal cord.

\medterm{Vertebrae} Vertebrae are the bones of the spinal column that protect the spinal cord, support the body, allow movement, and provide attachment for muscles and ligaments.

\medterm{Vertebral Artery} One of two arteries that travel through the neck and join to supply the brainstem, cerebellum, and back of the brain. Narrowing or blockage can cause dizziness, imbalance, or stroke.

\medterm{Vertebral Artery Dissection} A tear in the inner lining of a vertebral artery that lets blood enter the vessel wall, forming an intramural hematoma and possibly narrowing the artery or creating a pseudoaneurysm. It may occur spontaneously or after neck injury and can cause sudden neck or occipital headache, dizziness, imbalance, or a posterior-circulation transient ischemic attack or stroke.
\synonyms
VAD; Spontaneous Vertebral Artery Dissection

\medterm{Vertebral Canal} The passage formed by the stacked spinal bones that surrounds and protects the spinal cord. Narrowing, bleeding, tumors, fractures, or disc disease inside it can compress nerves and cause weakness or pain.

\medterm{Vertebral Column} The chain of bones running from the neck to the pelvis. It supports the head and trunk, protects the spinal cord, allows bending and twisting, and provides attachment for muscles and ribs.

\medterm{Vertebral Compression Fracture} Collapse or cracking of one of the bones in the spine, often from osteoporosis, injury, or cancer. It may cause sudden or gradual back pain, loss of height, or a rounded back and can sometimes press on nerves.

\medterm{Vertebral Hemangioma} Alternate name; see \textit{Spinal hemangioma}.

\medterm{Vertebral Tumor} A growth arising in or near a spinal bone. It may be benign or cancerous, start in the spine or spread there, and cause pain, fracture, nerve compression, weakness, or paralysis.

\medterm{Vertebrobasilar Circulatory Disorders} Conditions that reduce blood flow through the vertebral or basilar arteries supplying the brainstem, cerebellum, and posterior brain; symptoms may include vertigo, double vision, imbalance, weakness, or stroke.

\medterm{Vertebrobasilar Insufficiency} Reduced blood flow through arteries supplying the back of the brain. It may cause dizziness, double vision, imbalance, weakness, trouble speaking, or stroke-like symptoms and can be serious when sudden.

\medterm{Vertebrobasilar Insuffiency} Source spelling variant of \textit{vertebrobasilar insufficiency}, reduced blood flow in the vertebral or basilar circulation supplying the brainstem, cerebellum, and posterior brain.

\medterm{Vertebroplasty} A procedure in which medical cement is injected into a collapsed spinal bone to help stabilize it and sometimes reduce pain. It is used selectively because cement leakage, infection, nerve injury, or other complications can occur.

\medterm{Vertical} Oriented from superior to inferior or perpendicular to a horizontal plane. In medicine, it may describe position, transmission from parent to child, or a body-oriented imaging or movement axis.

\medterm{Vertical Talus} A stiff flatfoot present from birth in which one of the ankle bones points downward and the middle of the foot is bent upward. It may occur alone or with other conditions and usually needs early orthopedic treatment.

\medterm{Vertical Transmission} Passage of an infection or genetic condition from parent to child during pregnancy, delivery, or breastfeeding.

\medterm{Vertigo} A false feeling that you or your surroundings are moving, spinning, tilting, or shifting when they are not. It is a symptom rather than a disease and often comes from the inner ear, but brain and other conditions can also cause it.

\medterm{Vertigo-Associated Disorders} Vertigo-associated disorders cause an illusion of movement from vestibular, neurologic, visual, medication-related, or systemic dysfunction and may be accompanied by imbalance, nausea, or hearing symptoms.


\synonyms
Vestibular Disorders; Peripheral Vestibulopathy; Labyrinthine Disorders

\medterm{Very Long-Chain Acyl-CoA Dehydrogenase Deficiency} A rare inherited disorder in which cells cannot use certain long-chain fats for energy, especially during fasting or illness. Infants may develop low blood sugar, heart disease, or liver problems; older children or adults may have exercise-triggered muscle breakdown and dark urine.

\synonyms
VLCAD Deficiency; VLCADD

\medterm{Very-Low-Density Lipoprotein} A particle made by the liver that carries mainly triglycerides through the blood. High levels can contribute to fatty deposits in artery walls and increase cardiovascular risk, although the result is interpreted with the full lipid profile.

\synonyms
VLDL

\medterm{Vesical} Relating to the urinary bladder. For example, vesical inflammation affects the bladder lining, a vesical tumor is a bladder growth, and the vesical neck is the bladder outlet leading toward the urethra.

\medterm{Vesicant} A substance that causes blistering and tissue injury on contact. In oncology, vesicant drugs can cause severe extravasation injury if they escape from a vein into surrounding tissue.

\medterm{Vesicants} Substances that can cause severe blistering and tissue damage when they touch skin or escape from a vein. Some chemotherapy medicines and industrial chemicals are vesicants.

\medterm{Vesicle} A small, circumscribed, raised skin lesion containing clear, serous, or sometimes bloody fluid, usually less than 1 cm in diameter. A vesicle is a small blister; causes include viral infection, contact irritation, eczema, burns, friction, and immune-mediated skin disease.

\medterm{Vesicles} Plural of vesicle; see \textit{Vesicle}.

\medterm{Vesicoureteral Reflux} Backward flow of urine from the bladder toward the kidneys because the valve where a ureter enters the bladder does not work normally. It can cause repeated kidney infections and scarring, especially in children.

\synonyms
VUR; Urinary Reflux

\medterm{Vesico-Ureteric Reflux} Backward flow of urine from the bladder into one or both ureters and sometimes the kidneys, increasing the risk of recurrent urinary infection and renal scarring.

\medterm{Vesicovaginal and Ureterovaginal Fistula} Abnormal connections between the urinary tract and vagina: a vesicovaginal fistula connects bladder to vagina, while a ureterovaginal fistula connects ureter to vagina. Both commonly cause continuous urine leakage and require imaging and reconstructive urologic or gynecologic care.

\medterm{Vesicovaginal Fistula} An abnormal connection between the bladder and vagina that causes continuous or frequent leakage of urine through the vagina. It may follow childbirth injury, pelvic surgery, radiation, or advanced infection and usually requires specialist repair.

\medterm{Vesicular} Describing or containing small, fluid-filled blisters. A vesicular rash may occur with viral infection, allergy, eczema, burns, or other conditions; its appearance alone does not establish the cause.

\medterm{Vesicular Breathing} The normal soft, gentle breath sound heard over most outer lung areas with a stethoscope. Changes in loudness or quality may suggest fluid, collapse, blockage, infection, or other lung disease.

\medterm{Vesicular Palmoplantar Eczema} An itchy eczema affecting the palms or soles with small, deep-seated fluid-filled vesicles, followed by peeling, scaling, or fissures. Triggers include sweating, irritants, allergy, stress, and fungal infection elsewhere; persistent disease may need specialist treatment.

\medterm{Vesiculobullous Dermatoses} A group of skin diseases that cause small or large fluid-filled blisters. Causes include infection, allergy, autoimmune disease, inherited conditions, medicines, and injury; examination and sometimes a biopsy identify the cause.

\medterm{Vessel} A vessel is a tube carrying blood, lymph, air, urine, bile, or another fluid through the body, depending on context.

\medterm{Vestibular} Relating to the vestibular system of the inner ear and brain, which detects head motion and position and helps maintain balance and gaze stability.

\medterm{Vestibular Disorder} A problem affecting the inner-ear balance organs or the brain pathways that process balance signals. It may cause dizziness, vertigo, unsteadiness, nausea, abnormal eye movements, or difficulty walking. Causes include inner-ear disease, infection, head injury, medicines, migraine, reduced blood flow, or disorders of the brain or nerves.

\synonyms
Vestibular Balance Disorder

\medterm{Vestibular Migraine} Repeated attacks of dizziness or vertigo related to migraine, with or without headache. Episodes may include sensitivity to movement, light, sound, or visual patterns; other inner-ear and brain causes must be considered.

\medterm{Vestibular Neuritis} Sudden inflammation or dysfunction of the balance nerve, causing severe spinning, nausea, vomiting, and unsteadiness that may last for days. Hearing is usually unaffected, but urgent assessment may be needed to exclude a stroke.

\medterm{Vestibular Papillomatosis} A harmless normal variation in which several soft, small, finger-like bumps occur on the inner vulvar folds. It can resemble genital warts but is not caused by an infection and usually needs no treatment.

\medterm{Vestibular Rehabilitation} Exercise-based therapy that improves balance, gaze stability, and adaptation in people with
vestibular disorders or dizziness.

\medterm{Vestibular Schwannoma} A benign, slow-growing tumor of the vestibular portion of the eighth cranial nerve. See Acoustic Neuroma.

\medterm{Vestibular System} The inner-ear sensors and connected brain pathways that detect head movement and position and help maintain balance, posture, and steady vision. Problems can cause vertigo, unsteadiness, nausea, or blurred vision during movement.

\medterm{Vestibule} A balance area of the inner ear between the cochlea and semicircular canals. It contains sensors that detect head position and straight-line movement and send this information to the brain.

\medterm{Vestibulocochlear Nerve} The eighth cranial nerve, carrying information for hearing and balance from the inner ear to the brain.

\medterm{Vestibulodynia} Persistent burning, stinging, or tenderness at the entrance to the vagina, often triggered by touch, penetration, or inserting a tampon. Examination rules out infection and other causes; care may include pelvic-floor therapy, pain treatment, and support for sexual and emotional health.

\synonyms
Provoked Vestibulodynia; Vulvar Vestibulitis Syndrome; PVD

\medterm{Vestibulo-Ocular Reflex} An automatic response that moves the eyes opposite to head movement so vision stays steady while a person turns or walks. Problems with it can cause blurred vision, dizziness, or imbalance.

\medterm{Vestigial} Describing a structure or feature that has become greatly reduced through evolution and has little or no apparent original function.

\medterm{VEXAS Syndrome} A serious inflammatory illness that begins in adulthood because of a gene change acquired in blood-forming cells. It can cause persistent fever, painful joints, skin rashes, inflamed blood vessels or cartilage, anemia, and other blood abnormalities.

\medterm{V-Fib} Alternate name; see \textit{Ventricular fibrillation}.

\medterm{Viable} Capable of surviving or developing under suitable conditions. In clinical contexts it may describe a viable pregnancy, tissue, organism, cell, or treatment option.

\medterm{Vibration Problem} Symptoms or injury caused by repeated exposure to vibrating tools or equipment, which may affect nerves, blood vessels, muscles, and joints in the hands or arms.

\medterm{Vibrator} A device that produces mechanical vibration and may be used for sexual stimulation, therapeutic massage, or selected rehabilitation purposes.

\medterm{Vibrio} A genus of curved, Gram-negative bacteria; species include \textit{V. cholerae}, which causes cholera, and \textit{V. vulnificus}, which can cause severe wound or bloodstream infection.

\medterm{Vibrio Cholerae} A bacterium that causes cholera, a severe diarrheal illness spread mainly through contaminated water or food. It can cause sudden, large-volume watery diarrhea, dehydration, shock, and death without rapid fluid replacement.

\medterm{Vibrio Infection} Infection caused by Vibrio bacteria, acquired through contaminated seafood or exposure of wounds to warm seawater. It may cause watery diarrhea, wound infection, or rapidly worsening bloodstream infection, especially in people with liver disease or weakened immunity. Severe wound pain, blistering, fever, or shock can occur.

\synonyms
Vibriosis; Vibrio Bacteremia

\medterm{Vibrio Parahaemolyticus Infection} An acute halophilic bacterial gastroenteritis acquired through ingestion of raw or undercooked marine seafood, particularly shellfish. It typically manifests with explosive watery diarrhea, abdominal cramps, nausea, vomiting, low-grade fever, and headache within 12 to 24 hours of exposure, usually resolving spontaneously within 3 to 5 days without antibiotics.

\synonyms
Vibrio Parahaemolyticus Gastroenteritis; Halophilic Vibrio Infection

\medterm{Vibrios} Comma-shaped bacteria, usually Gram-negative, that include species capable of causing cholera and wound or gastrointestinal infections.

\medterm{Vibrio vulnificus Infection} Infection caused by a saltwater-associated bacterium that can produce rapidly progressive wound infection, sepsis, or severe gastroenteritis. Risk is higher with liver disease, iron overload, or immune suppression; raw oysters and seawater exposure are important sources.

\medterm{Vibrio Vulnificus Infections} Infections caused by Vibrio vulnificus, a gram-negative bacterium found in warm seawater and raw
seafood. It can cause wound infection, bloodstream infection, or gastroenteritis, with
particularly severe disease in people with chronic liver disease or other risk factors.

\medterm{Video-Assisted Thoracoscopic Surgery} A chest operation performed through several small openings using a camera and long instruments rather than a large incision between the ribs. It may be used for lung or pleural biopsy, removing part of a lung, or other chest procedures and often allows faster recovery, though bleeding, air leakage, infection, or conversion to open surgery can occur.

\synonyms
VATS; Video Thoracoscopy; Thoracoscopic Surgery

\medterm{Videofluoroscopic Swallow Study} A moving X-ray test that records swallowing after small amounts of food or liquid containing contrast are taken. It shows whether food enters the airway, where swallowing is delayed, and which textures or techniques may be safer.

\medterm{Video Laryngoscope} An airway device with a small camera that shows the entrance to the windpipe on a screen. It helps clinicians place a breathing tube, especially when a direct view is difficult, but does not remove the risks of the procedure.

\medterm{Video Laryngoscopy} A technique for viewing the entrance to the windpipe with a camera on a blade connected to a screen. It can help clinicians place a breathing tube when a direct view is difficult.

\medterm{Videostroboscopy} A voice examination that uses a camera and flashing light to create a slow-motion view of the vocal cords as they vibrate. It helps investigate hoarseness, voice strain, and vocal-cord injury.

\medterm{Vildagliptin} An oral dipeptidyl-peptidase-4 inhibitor used in some countries for type 2 diabetes, alone or with other medicines. It improves glucose-dependent insulin secretion; liver function and pancreatitis symptoms may require attention.

\medterm{Villi} Tiny finger-like projections lining the small intestine that greatly increase the area available to absorb nutrients. Similar projections in the placenta help exchange substances between the pregnant person and fetus.

\medterm{Villous Adenoma} A high-risk subtype of neoplastic colonic adenomatous polyp characterized histologically by long, finger-like, frond-like epithelial projections covering more than 75% of the polyp mass, associated with a high potential for malignant transformation and excessive secretory mucus/electrolyte loss (secretory diarrhea with hypokalemia).
\synonyms
Villous Colonic Polyp; Papillary Adenoma

\medterm{Villous Atrophy} Flattening or loss of the tiny projections lining the small intestine. This reduces nutrient absorption and may occur in celiac disease, some infections, immune disorders, or other intestinal conditions.

\medterm{Villus} A small finger-like projection, especially one of the intestinal mucosa that increases absorptive surface area; chorionic villi mediate exchange between placenta and fetus.

\medterm{Villus Atrophy} The histopathological flattening, blunting, and loss of the microscopic finger-like villi lining the small intestinal mucosa, accompanied by crypt hyperplasia and intraepithelial lymphocytosis, classically observed in celiac disease, tropical sprue, and autoimmune enteropathy.

\medterm{Vinca Alkaloids} Cancer medicines including vincristine and vinblastine that disrupt microtubules, structures needed for cell division. They can treat several cancers but may cause nerve damage, low blood counts, constipation, or other toxicity.

\medterm{Vincent Angina} An acute necrotizing infection of the gums and oropharynx associated with fusobacteria and spirochetes. It causes painful ulceration, bleeding, foul breath, and gingival necrosis and requires dental or medical treatment.

\medterm{Vincent Disease} Acute necrotizing ulcerative gingivitis associated with fusospirochetal bacteria, causing painful bleeding gums, ulceration, and foul breath.

\medterm{Vincent Infection} Vincent infection is acute necrotizing ulcerative gingivitis caused by mixed oral bacteria, producing painful bleeding gums, ulcers, and foul breath.

\medterm{Vincent's Angina} Acute necrotizing ulcerative gingivitis or tonsillitis associated with fusiform bacteria and spirochetes, causing painful ulcers, bleeding, and fetid breath.

\medterm{Vincent Stomatitis} A painful necrotizing infection of the gums caused by mixed anaerobic bacteria, producing foul breath, bleeding, ulceration, and gray necrotic tissue; it is also called acute necrotizing ulcerative gingivitis.

\medterm{Vinorelbine} A chemotherapy medicine that interferes with cell division. It is used for selected lung and breast cancers and can cause low blood counts, infection, constipation, nerve symptoms, nausea, or irritation if it leaks outside a vein.

\medterm{Viomycin} An older aminoglycoside antibiotic related to capreomycin that was used against tuberculosis and other mycobacterial infections; kidney and hearing toxicity limit its use.

\medterm{VIPoma} A rare functioning neuroendocrine tumor arising from the non-beta islet cells of the pancreas (or occasionally extra-pancreatic sympathetic ganglia in children) that hypersecretes vasoactive intestinal peptide (VIP), causing refractory secretory diarrhea and hypokalemia.
\synonyms
Vasoactive Intestinal Peptide Tumor; Verner-Morrison Neoplasm

\medterm{VIPoma Syndrome} A rare neuroendocrine tumor syndrome arising predominantly from the non-beta islet cells of the pancreas that hypersecretes vasoactive intestinal peptide (VIP). Clinically characterized by massive, secretory, watery diarrhea persisting during fasting, severe hypokalemia, and hypochlorhydria or achlorhydria (the WDHA syndrome), accompanied by dehydration and metabolic acidosis.

\synonyms
WDHA Syndrome; Verner-Morrison Syndrome; Pancreatic Cholera

\medterm{Viraemia} The presence of viruses in the bloodstream. It may be brief or long-lasting and can allow the infection to spread from its original site to other tissues.

\medterm{Viral} Caused by or relating to a virus. Viruses enter living cells to multiply and can cause illnesses ranging from mild colds to serious brain, lung, liver, or bleeding diseases; testing and treatment depend on the specific virus.

\medterm{Viral Antibody} A protective protein made by the immune system after infection or vaccination that recognizes a virus. Blood tests for these antibodies may show past exposure, immune response, or selected current infections, depending on timing and test design.

\medterm{Viral Arthritis} Joint inflammation caused by or associated with a viral infection. It may cause sudden pain, swelling, and stiffness in several joints and often improves as the infection resolves, although some viruses can cause prolonged symptoms.

\medterm{Viral Assembly} The stage of viral reproduction when newly made viral genetic material and proteins are put together to form complete new virus particles.

\medterm{Viral Attachment} The first step of infection, when proteins on a virus bind to matching molecules on a host cell. This binding helps determine which species, tissues, or cell types the virus can infect.

\medterm{Viral Carcinogenesis} Cancer development caused or promoted by a viral infection. Some viruses alter cell growth directly, weaken immune protection, or cause long-lasting inflammation; examples include human papillomavirus, hepatitis B and C viruses, and Epstein-Barr virus.

\medterm{Viral Conjunctivitis} Infection of the thin lining over the eye and inside the eyelids, usually causing redness, watering, irritation, and discharge. It spreads easily through contact with infected secretions or contaminated hands and surfaces.

\medterm{Viral Culture} Growth of a virus in susceptible cell cultures to identify an infectious agent. It can be useful for
selected viruses, but molecular and antigen-based tests are often faster or more sensitive;
the appropriate method depends on the suspected infection and specimen.

\medterm{Viral Encephalitis} Inflammation of the brain caused by a virus. It may cause fever, headache, confusion, behavior change, seizures, weakness, or reduced consciousness. Severe cases can be life-threatening.

\medterm{Viral Entry} The process by which a virus gets inside a host cell and releases its genetic material so the cell can make more virus.

\medterm{Viral Exanthem} A widespread rash caused by a viral infection. It may occur with fever, tiredness, sore throat, cough, diarrhea, or other symptoms; the rash pattern and overall illness help determine the cause.

\medterm{Viral Gastroenteritis} Alternate name; see \textit{Stomach Flu}.

\medterm{Viral Haemorrhagic Fever} A severe illness caused by certain viruses that can damage blood vessels, disturb clotting, and cause bleeding, shock, or organ failure. The way it spreads and its treatment depend on the virus.

\medterm{Viral Hemorrhagic Fevers} Severe illnesses caused by certain viruses that can damage blood vessels, disrupt clotting, and lead to bleeding, shock, or organ failure. Some spread easily and may require special infection-control measures.
\synonyms
VHF; Hemorrhagic Fevers; Viral Hemorrhagic Disease

\medterm{Viral Hepatitis} Liver inflammation caused by hepatitis viruses A, B, C, D, or E. Hepatitis A and E usually spread through contaminated food or water; B, C, and D can spread through blood or body fluids and may become long-term infections that damage the liver.

\synonyms
Infectious Hepatitis

\medterm{Viral Hepatitis Types} Five viruses that can inflame the liver. They differ in how they spread, whether infection can become long-term, the complications they cause, and whether vaccination or curative treatment is available.

\medterm{Viral Immune Evasion} Ways viruses avoid or weaken the body’s immune defenses, allowing them to survive, multiply, or return after infection. These methods can make infections harder to clear and vaccines or treatments less effective.

\medterm{Viral Infection} An illness caused by a virus entering body cells and making copies of itself. Viruses can damage cells directly or cause inflammation as the immune system responds. Viral infections range from mild illnesses such as colds and cold sores to serious diseases such as influenza, COVID-19, hepatitis, and HIV. Some viruses remain inactive in the body and can become active again later, such as the chickenpox virus causing shingles.

\synonyms
Viral Infections; Viral Disease; Viral Illness

\medterm{Viral Infections and Pregnancy} Viral infections during pregnancy may affect maternal illness, fetal development, placenta, or newborn health. Evaluation considers the virus, gestational age, symptoms, immunity, exposure, testing, treatment, vaccination, and measures to reduce perinatal transmission.

\medterm{Viral Infections of the Mouth} Infections of oral skin or mucosa caused by viruses such as herpes simplex, varicella-zoster, enteroviruses, HPV, or EBV. They may cause blisters, ulcers, warts, sore throat, or systemic symptoms, and treatment depends on the virus and host risk.

\medterm{Viral Latency} A state in which a virus remains inside cells without actively making new virus particles. Some latent viruses can become active again, especially when the immune system is weakened or the body is under stress.

\medterm{Viral Latency Mechanisms} The ways some viruses remain inactive inside cells while avoiding immune detection. Reactivation can occur when immune defenses are weakened or after other triggers, depending on the virus.

\medterm{Viral Load} The amount of virus or viral genetic material in a body fluid or tissue, measured by a laboratory assay; trends can help assess disease activity or response to antiviral treatment.

\medterm{Viral Meningitis} Infection and inflammation of the coverings of the brain and spinal cord caused by a virus. It may cause fever, headache, neck stiffness, vomiting, light sensitivity, or confusion. Bacterial meningitis can cause similar symptoms and is more dangerous.

\medterm{Viral Myocarditis} Inflammation of the heart muscle caused by a viral infection or by the immune response after an infection. It may cause tiredness, chest discomfort, palpitations, heart failure, shock, or dangerous abnormal heart rhythms.

\synonyms
Infectious Viral Myocarditis; Acute Viral Myocarditis

\medterm{Viral Pathogenesis} The ways a virus enters the body, infects and spreads through cells, and causes illness. Damage may come directly from infected cells or from the immune system’s response to the infection.

\medterm{Viral Pericarditis} Inflammation of the sac around the heart caused by a viral infection. It often follows a respiratory or stomach illness and may cause sharp chest pain that changes with breathing or position.

\synonyms
Acute Viral Pericarditis; Viral Pericardial Inflammation

\medterm{Viral Pharyngitis} Pharyngeal inflammation caused by viral pathogens such as adenovirus, rhinovirus, coronavirus, Epstein--Barr virus, or influenza virus. Characterized by cough, rhinorrhea, conjunctivitis, and lack of prominent tonsillar exudates, resolving with supportive care.

\synonyms
Viral Sore Throat; Acute Viral Pharyngotonsillitis; Viral Upper Respiratory Pharyngitis

\medterm{Viral Pneumonia} Infection and inflammation of lung tissue caused by a virus, producing cough, fever, chills, breathlessness, chest discomfort, fatigue, and variable impairment of oxygenation. Severe disease can cause respiratory failure or secondary bacterial infection.

\medterm{Viral Release} The stage when newly made viruses leave an infected cell and spread to other cells. Some leave by pushing through the cell membrane, while others are released when the infected cell breaks apart.

\medterm{Viral Replication} The process by which a virus uses an infected cell to copy its genetic material and make viral proteins. The exact steps differ among viruses, but the result is production of new virus particles.

\medterm{Viral Replication Cycle} The steps a virus uses to make more copies inside a living cell. Usually the virus attaches to the cell, enters, releases its genetic material, makes new parts, assembles new viruses, and leaves to infect other cells.

\medterm{Viral Rhinitis} Acute inflammation and edema of the nasal mucous membranes caused by respiratory viral pathogens, predominantly rhinoviruses, coronaviruses, respiratory syncytial virus, or parainfluenza viruses. Clinically recognized as acute viral coryza or the common cold, it manifests with clear anterior rhinorrhea, nasal congestion, sneezing, itchy throat, low-grade malaise, and temporary hyposmia, generally resolving without antimicrobial therapy within 7 to 10 days.

\synonyms
Acute Viral Rhinitis; Acute Coryza; Common Cold Rhinitis

\medterm{Viral Structure} The parts of a virus: genetic material made of DNA or RNA inside a protective protein coat. Some viruses also have an outer fatty covering taken from the infected cell’s membrane.

\medterm{Viral Uncoating} The step after a virus enters a cell when its protective outer coat comes apart and releases the genetic material. The cell can then be directed to make new viral components.

\medterm{Virchow Node} A firm, non-tender, pathologically enlarged left supraclavicular lymph node (Troisier sign) located in the left supraclavicular fossa, which receives abdominal lymphatic drainage via the thoracic duct and signals metastatic spread of abdominal malignancies (particularly gastric adenocarcinoma).
\synonyms
Troisier Node; Left Supraclavicular Sentinel Node

\medterm{Virchow-Robin Spaces} Small fluid-filled spaces around blood vessels as they pass through the brain. They are a normal part of the spaces that surround and support these vessels.

\synonyms
Perivascular Spaces

\medterm{Virchow's Triad} Three factors that increase the chance of a blood clot forming: slow blood flow, injury to the inside of a blood vessel, and blood that clots too easily. These factors can contribute to clots in deep veins or clots that travel to the lungs.
\synonyms{Virchow Triad; Triad of Virchow; Thrombogenesis Triad}

\medterm{Virchow Triad} Three factors that increase the chance of a blood clot forming: injury to the inside of a blood vessel, slow blood flow, and blood that clots too easily.
\synonyms
Virchow's Triad; Thrombosis Pathophysiological Triad

\medterm{Viremia} The presence of a virus in the bloodstream. It may be brief or persistent and can allow infection to spread to other organs. Laboratory tests may detect the virus itself or its genetic material.

\medterm{Viridans Streptococci} Viridans streptococci are a group of bacteria commonly found in the mouth that can cause dental infections and, after entering the blood, infect heart valves.

\medterm{Virilisation} The development of male secondary sexual characteristics in females, resulting from
excessive androgen production or exposure.

\medterm{Virilism} Development of masculinizing features in a female, usually from excess androgen exposure or production, such as hirsutism, deep voice, acne, and menstrual disturbance.

\medterm{Virilism Foetal} Prenatal or neonatal masculinization of a genetically female fetus or infant due to
excess androgen exposure. Causes include congenital adrenal hyperplasia, maternal
androgen-producing disease, or placental and fetal steroidogenic disorders.

\medterm{Virilization} Virilization is development of typically male secondary sexual features, such as hirsutism, deep voice, acne, increased muscle mass, or clitoromegaly, from androgen excess or androgen exposure.

\medterm{Virilize} To virilize is to develop or cause typically male physical changes from androgen excess or treatment, including facial hair, voice deepening, acne, muscle change, or genital enlargement.

\medterm{Virology} The study of viruses, including their structure, reproduction, spread, effects on the body, diagnosis, treatment, and prevention.

\medterm{Virotherapy} Treatment that uses a virus, often modified, to attack cancer cells or deliver a useful gene into cells. It is an evolving field; the possible benefits, risks, and availability depend on the specific therapy and disease.

\medterm{Virtual Colonoscopy} Virtual colonoscopy is CT colonography, which uses bowel preparation and gas or carbon-dioxide insufflation to create images of the colon; it can detect polyps and masses but cannot remove lesions or biopsy them.

\medterm{Virulence} How severely a microorganism can cause disease or tissue damage. It depends on the organism’s harmful properties and on the person’s immune defenses and other health factors.

\medterm{Virulence Factor} A feature of a microorganism that helps it survive in the body, avoid immune defenses, attach to cells, damage tissue, or cause disease. Examples include toxins, protective coatings, and enzymes.

\medterm{Virulent} Capable of causing severe disease or extensive harm. Virulence describes the harmfulness of a microorganism and is different from infectivity, which describes how easily it spreads or establishes infection.

\medterm{Virus} An infectious agent made of genetic material inside a protective protein coat, sometimes with an outer membrane. It can reproduce only inside living cells and may cause illness by damaging cells or triggering inflammation.

\medterm{Virus Activation} Virus activation is the switch from a dormant or inactive viral state to production of new viral components and particles.

\medterm{Virus Assembly} Virus assembly is the stage of the viral life cycle when viral genetic material and proteins are packaged into new virus particles.

\medterm{Virus Diseases} Illnesses caused by viruses. They range from mild infections of one area, such as a cold, to severe diseases affecting the lungs, liver, brain, blood, or several organs; prevention and treatment depend on the virus.

\medterm{Viruses} Viruses are infectious particles containing DNA or RNA within a protein coat, sometimes with a lipid envelope; they replicate only inside host cells and cause disease through direct injury or immune responses.

\medterm{Virus Integration} Virus integration is insertion of viral genetic material into the DNA of a host cell, a process used by some viruses and important in certain cancers.

\medterm{Virus Protein} A virus protein is a protein made from viral genetic instructions that helps form the virus, enter cells, replicate, or evade immunity.

\medterm{Virus Receptor} A virus receptor is a molecule on a host cell that a virus binds to in order to attach to and enter the cell.

\medterm{Virus Replication} Virus replication is the process by which a virus copies its genetic material and makes proteins and new virus particles inside a host cell.

\medterm{Virus Uncoating} Virus uncoating is removal of the virus's protective outer layer after entry into a host cell, releasing its genetic material for replication.

\medterm{Viscera} The internal organs of the body, particularly those within the thoracic and abdominal
cavities. The thoracic viscera include the heart, lungs, esophagus, trachea, and thymus.

\medterm{Visceral} Relating to the internal organs, especially those within the thorax or abdomen, rather than the body wall or limbs.

\medterm{Visceral Afferents} Sensory nerve fibers that carry information from internal organs to the spinal cord and brain. They convey signals related to stretch, chemical changes, inflammation, and internal pain and help regulate autonomic and reflex responses.

\medterm{Visceral Larva Migrans} A syndrome caused by migration of animal nematode larvae through human tissues, most often from Toxocara species. It may produce fever, eosinophilia, liver lesions, or ocular and neurologic disease.

\medterm{Visceral Leishmaniasis} A severe, potentially fatal systemic protozoan infection caused by species of the Leishmania donovani complex (Leishmania donovani and Leishmania infantum) transmitted to humans by the bite of female phlebotomine sandflies. Characterized by the intracellular dissemination of amastigotes within the reticuloendothelial system, it presents with prolonged irregular fever, progressive massive splenomegaly, hepatomegaly, profound pancytopenia (anemia, leukopenia, and thrombocytopenia), hypergammaglobulinemia, and severe wasting, fatal in over 95\% of untreated cases.

\synonyms
Kala-Azar; Black Fever; Dumdum Fever; Systemic Leishmaniasis

\medterm{Visceromegaly} Abnormal enlargement of one or more internal organs, such as the liver, spleen, or kidneys. It is a physical finding with causes that include infection, congestion, storage disease, inflammation, and cancer.

\medterm{Viscus} An internal organ, especially one in the chest or abdomen. The plural is viscera; examples include the heart, lungs, liver, stomach, and intestines.

\medterm{Vision} The sensory ability to perceive light and interpret visual information processed by the eyes and brain, enabling recognition of shapes, colors, depth, and motion. Normal vision depends on the proper function of the cornea, lens, retina, and optic nerve, as well as intact neural pathways to the visual cortex.

\synonyms
Eyesight; Sight; Visual Perception

\medterm{Vision And Hearing} Two major senses used to understand and communicate with the world. Problems may develop separately or together because of aging, illness, injury, medicines, or inherited conditions and can affect learning, safety, communication, and daily activities.

\medterm{Vision Impairment} Reduced visual acuity or visual-field function that affects daily activities. Causes include refractive
error, cataract, glaucoma, retinal disease, optic-nerve injury, and neurologic disease.
Evaluation and treatment depend on the cause; rehabilitation and visual aids may improve
function.

\medterm{Vision Loss} A decrease in the ability to see that cannot be fully corrected with eyeglasses, contact
lenses, or medical treatment. It ranges from partial impairment to complete blindness
and can result from conditions such as macular degeneration, glaucoma, diabetic
retinopathy, cataracts, and traumatic injury.

\medterm{Vision Problems} Changes in visual acuity, field, color, or clarity caused by refractive error, eye
disease, neurologic disease, or systemic illness. Sudden vision loss is an emergency.

\medterm{Vision Receptors} Light-sensing cells in the retina, mainly rods and cones, that convert light into nerve signals. Rods support dim-light vision, while cones provide color and sharp detail in brighter conditions.

\medterm{Vision Screening} Brief checks used to find possible eye or vision problems before symptoms are obvious. Newborn checks may detect cataracts or other serious disorders; screening in young children looks for poor vision, misaligned eyes, or risks for reduced vision and leads to a full eye examination when needed.

\medterm{Vision Test} An examination that measures visual acuity, fields, color, alignment, or eye structures to detect refractive error, retinal disease, optic-nerve dysfunction, glaucoma, or other causes of visual impairment.

\medterm{Visual Acuity} How clearly a person can see fine details, usually measured by reading letters or symbols at a set distance. It can be reduced by focusing problems, eye disease, or disorders affecting the visual nerves or brain.

\medterm{Visual Acuity Test} A test measuring the smallest standardized detail a person can resolve at a specified distance, usually reported as a fraction or decimal.

\medterm{Visual Agnosia} Inability to recognize familiar objects, faces, or symbols by sight even though the eyes may see adequately. It usually results from damage to brain areas that interpret vision, such as after a stroke or head injury.

\medterm{Visual Evoked Potential} A test that records the brain’s electrical response to flashes or changing visual patterns. It helps assess the nerve pathways that carry visual signals from the eyes to the brain.
\synonyms
VEP; Visual Evoked Response (VER)

\medterm{Visual Field} The entire area that can be seen while the eyes remain fixed in one position. Visual-field loss may result from retinal, optic-nerve, or brain lesions and is assessed by perimetry.

\medterm{Visual Field Defect} Loss or change in part of the area a person can see while looking straight ahead. It may result from disease of the retina, optic nerve, or brain, and the pattern can help identify where the problem is.

\medterm{Visual Field Perimetry} A test that maps how well a person sees in different parts of the visual field while looking at a fixed point. It can detect and monitor missing areas caused by eye-nerve, retinal, or brain disease.

\medterm{Visual Fields} The total area that can be seen while the eyes are fixed in one position; defects may indicate retinal, optic-nerve, or brain disease.

\medterm{Visual Field Test} A test that measures central and side vision and maps areas that are missing or less sensitive.

\medterm{Visual Impairment} Reduced vision that cannot be fully corrected by ordinary glasses or contact lenses. Causes include refractive error, cataract, glaucoma, retinal disease, optic-nerve injury, and disorders of the brain or nervous system.

\medterm{Visualized Laser-Assisted Prostatectomy} A procedure that uses a laser, guided by a camera, to remove or vaporize prostate tissue blocking urine flow. It is used for some people with an enlarged prostate and may cause bleeding, infection, scar narrowing, or temporary urinary and sexual problems.

\synonyms
VLAP

\medterm{Visual Neglect} A brain disorder in which a person fails to notice or respond to things on one side, even though the eyes may work normally. It often follows a stroke affecting the right side of the brain and causes neglect of the left side.

\synonyms
Hemispatial Neglect; Spatial Inattention

\medterm{Vital} Essential for life or relating to vital physiologic functions. Vital signs commonly include temperature, pulse, respiratory rate, blood pressure, and oxygen saturation.

\medterm{Vital Capacity} The greatest amount of air a person can breathe out after filling the lungs as fully as possible. It is measured during lung-function testing and may be reduced by lung disease, weak breathing muscles, or restriction of chest movement.

\medterm{Vital Signs} Basic measurements of physiologic status, commonly including temperature, pulse, respiratory rate, blood
pressure, and oxygen saturation. Values vary with age, activity, illness, and measurement
conditions; trends and clinical context are important.

\medterm{Vitamin A} A fat-soluble vitamin needed for vision, immune function, skin and mucosal health, and growth; deficiency causes night blindness, while excess can injure the liver and developing fetus.
\synonyms
Retinol; Retinoic Acid

\medterm{Vitamin A Blood Test} A blood test measuring circulating retinol, used to assess suspected vitamin A deficiency or toxicity in appropriate clinical settings.

\medterm{Vitamin A Deficiency} A nutritional disorder resulting from inadequate intake or absorption of vitamin A (retinol), causing dry eyes (xerophthalmia), night blindness, dry skin, and weakened immunity. Severe deficiency leads to corneal ulceration and is the leading cause of preventable childhood blindness worldwide; clinical eye stages are categorized using the WHO xerophthalmia classification.

\synonyms
Hypovitaminosis A; Retinol Deficiency; Vitamin A Lack

\medterm{Vitamin A: Retinal} Vitamin A (retinol) is a fat-soluble vitamin essential for vision, immune function, cell
differentiation, and reproduction. In the eye, retinol is oxidized to retinal
(retinaldehyde), which binds to opsin in the photoreceptor outer segments forming
rhodopsin (in rods) and photopsins (in cones).

\medterm{Vitamin A Toxicity} Harm caused by taking too much vitamin A or related medicines. It may cause nausea, vomiting, headache, dry skin, hair loss, liver injury, high calcium levels, and serious fetal abnormalities during pregnancy.

\medterm{Vitamin B} Vitamin B refers to a group of water-soluble vitamins that support energy metabolism, blood-cell formation, nerve function, and other processes. The group includes several distinct vitamins, not one single substance.

\medterm{Vitamin B1} Thiamine, a water-soluble vitamin that acts as a cofactor in carbohydrate and amino-acid metabolism and supports nerve function. Severe deficiency can cause beriberi or Wernicke encephalopathy.
\synonyms
Thiamine; Aneurin

\medterm{Vitamin B12} A vitamin needed for red-cell formation, DNA production, and normal nerve function. Deficiency can cause anemia, numbness, balance problems, or cognitive changes.
\synonyms
Cobalamin; Cyanocobalamin

\medterm{Vitamin B12: Cobalamin} Alternate name; see \textit{Vitamin B12}.

\medterm{Vitamin B12 Deficiency} Too little vitamin B12, which can cause large abnormal red cells, tiredness, numbness, balance problems, or nerve damage. Causes include poor intake, poor absorption, autoimmune disease, stomach surgery, and medicines.

\medterm{Vitamin B12 Deficiency Anemia} Anemia caused by too little vitamin B12 or poor absorption of it. It can cause tiredness, pale skin, a sore tongue, numbness, balance problems, and permanent nerve damage if untreated for a long time.

\synonyms
Vitamin B12 Deficiency Anaemia

\medterm{Vitamin B12 Deficiency In Older Adults} Deficiency caused by inadequate intake, malabsorption, autoimmune intrinsic-factor
deficiency, gastric surgery, or medications. It may cause megaloblastic anemia,
neuropathy, cognitive changes, gait impairment, and glossitis; neurologic injury can
occur without anemia.

\medterm{Vitamin B12 Level} The amount of vitamin B12 measured in blood; results are interpreted with symptoms and related tests because serum levels may not fully reflect tissue deficiency.

\medterm{Vitamin B1: Thiamine} Alternate name; see \textit{Vitamin B1}.

\medterm{Vitamin B2} Riboflavin, a water-soluble vitamin that helps cells produce energy and supports healthy skin, eyes, and the lining of the mouth. Deficiency can cause cracks at the mouth corners, a sore tongue, skin changes, or anemia.
\synonyms
Riboflavin; Lactoflavin

\medterm{Vitamin B2 Riboflavin} Alternate name; see \textit{Riboflavin}.

\medterm{Vitamin B3} A vitamin that helps cells use food for energy and supports the skin and nervous system. Severe deficiency causes pellagra, with diarrhea, skin inflammation, confusion, and other neurologic problems.

\synonyms
Nicotinic Acid; Nicotinamide

\medterm{Vitamin B5} A water-soluble vitamin needed to help the body use fats, carbohydrates, and proteins and make hormones and other substances. It is found in many foods, so deficiency is rare; very large supplement doses may cause stomach upset.

\synonyms
Pantothenate; Calcium Pantothenate

\medterm{Vitamin B6} A group of related vitamins needed for amino-acid metabolism, red-cell production, and nerve function. Excess supplements can cause sensory nerve damage.
\synonyms
Pyridoxine; Pyridoxal Phosphate

\medterm{Vitamin B6: Pyridoxine} Alternate name; see \textit{Vitamin B6}.

\medterm{Vitamin B7} A vitamin that helps the body use fats, carbohydrates, and proteins. Deficiency is uncommon but may cause hair loss, skin rash, or neurologic symptoms; large doses can interfere with some laboratory tests.
\synonyms
Vitamin H; Coenzyme R

\medterm{Vitamin B9} A vitamin needed for DNA production and healthy red blood cells. Deficiency can cause large-cell anemia and, during pregnancy, increases the risk of neural-tube defects; folic acid is recommended before conception and early in pregnancy.

\synonyms
Folic Acid; Pteroylglutamic Acid

\medterm{Vitamin B Complex} A group of eight water-soluble vitamins that act mainly as coenzymes or their precursors in cell metabolism. The group includes thiamine (B1), riboflavin (B2), niacin (B3), pantothenic acid (B5), pyridoxine (B6), biotin (B7), folate (B9), and cobalamin (B12), which together support energy metabolism, red-blood-cell formation, and nervous-system function. These vitamins occur together in whole grains, meat, eggs, and legumes; individual deficiencies cause distinct disorders, and combined supplements can prevent or correct deficiency. The sequence has gaps because early candidate substances were later excluded as nonessential: B8 (inositol) is synthesized by the body, B4 (choline, adenine) is not a true vitamin, and the provisional labels B10 and B11 were abandoned once their activity was identified as folate or para-aminobenzoic acid.
\synonyms
B Vitamins; Vitamin B Group

\medterm{Vitamin C} A water-soluble vitamin needed for collagen formation, wound healing, immune function, antioxidant activity, and absorption of iron from plant foods. Severe deficiency causes scurvy, with bleeding gums, bruising, poor wound healing, and weakness.
\synonyms
Ascorbic Acid

\medterm{Vitamin C Deficiency} A lack of vitamin C, a water-soluble vitamin needed for
collagen formation, wound healing, iron absorption, and antioxidant activity. It follows
prolonged low intake and, when severe, causes scurvy, with bleeding gums, easy bruising,
poor wound healing, joint pain, and weakness. Replacement and dietary correction reverse
the deficiency.

\medterm{Vitamin D} Vitamin D supports calcium and phosphate balance, bone mineralization, muscle function, and immune processes; deficiency causes rickets in children and osteomalacia or bone loss in adults.

\medterm{Vitamin D3} Cholecalciferol, a form of vitamin D made in skin and found in some foods and supplements. It helps the body absorb calcium and maintain healthy bones.

\medterm{Vitamin D: Calcitriol} Alternate name; see \textit{Calcitriol}.

\medterm{Vitamin D Deficiency} Too little vitamin D for the body’s needs. It can result from low intake or sunlight, poor absorption, or liver or kidney disease and may weaken bones, cause muscle weakness, or lead to rickets in children.

\medterm{Vitamin Deficiency Anemia} Anemia caused by inadequate vitamin B-12, folate, or other nutrients needed for normal red blood cell production.

\medterm{Vitamin E} Vitamin E is a group of fat-soluble compounds with antioxidant activity. It is found in foods such as nuts, seeds, and vegetable oils; high-dose supplements can increase bleeding risk and interact with medicines.
\synonyms
Tocopherol; Alpha-Tocopherol

\medterm{Vitamin E Deficiency} Deficiency of a fat-soluble antioxidant vitamin, usually caused by fat-malabsorption disorders or rare inherited conditions. Severe deficiency can cause peripheral neuropathy, ataxia, muscle weakness, retinal dysfunction, or hemolytic anemia.

\medterm{Vitamin E: Tocopherol} Alternate name; see \textit{Vitamin E}.

\medterm{Vitamin K} A fat-soluble vitamin required to activate several blood-clotting proteins and support bone metabolism; the main dietary form is phylloquinone, found especially in leafy green vegetables and some vegetable oils. Deficiency can cause bleeding and is treated or prevented with replacement.

\medterm{Vitamin K Deficiency} Too little vitamin K for normal production of clotting factors, causing an increased bleeding tendency. It can result from poor intake, malabsorption, liver disease, or certain medicines; newborns receive vitamin K prophylaxis.

\medterm{Vitamin K Deficiency Bleeding} Bleeding in a newborn or young infant caused by inadequate vitamin K-dependent clotting activity. It may occur early, classically, or late and can involve the brain or gastrointestinal tract.

\synonyms
VKDB; Neonatal Hypoprothrombinemia

\medterm{Vitamin K: Phylloquinone} Alternate name; see \textit{Vitamin K}.

\medterm{Vitamin K Prophylaxis} Giving vitamin K to a newborn soon after birth to prevent bleeding caused by vitamin K deficiency.

\medterm{Vitamins} Organic nutrients needed in small amounts for normal metabolism, growth, and tissue maintenance; some are water-soluble and others fat-soluble, and both deficiency and excess can cause disease.

\medterm{Vitamin Toxicity} Harm from excessive vitamin intake, most often involving fat-soluble vitamins that accumulate in the body. Vitamin A can cause headache, liver injury, and bone effects; vitamin D can cause hypercalcemia; and high-dose vitamin E can increase bleeding risk. Stop the source and treat the affected organ or metabolic abnormality.

\medterm{Vitiligo} Vitiligo is a condition in which areas of skin lose pigment because melanocytes are destroyed or stop functioning. It is not contagious, and treatment may include camouflage, topical medicines, or light therapy.

\medterm{Vitrectomy} Eye surgery that removes some or all of the vitreous, the clear gel that fills the middle of the eye. Small instruments are inserted through the wall of the eye to remove cloudy gel, blood, scar tissue, or a foreign object and to repair problems of the retina. It may be used for retinal detachment, diabetic eye disease, macular hole, macular pucker, severe eye injury, or serious eye infection. The removed gel may be replaced with a salt solution, air or gas, or silicone oil.

\medterm{Vitreous} Relating to the transparent gel filling most of the space between the lens and retina, or describing a glass-like appearance. Vitreous disorders can cause floaters, hemorrhage, detachment, or retinal traction.

\medterm{Vitreous Body} The transparent gel filling the large space between the lens and retina. It helps maintain the shape of the eye and transmits light; age-related liquefaction can cause floaters or contribute to retinal traction.

\synonyms
Vitreous Humor; Corpus Vitreum

\medterm{Vitreous Detachment} Separation of the vitreous gel from the retina, usually related to aging and often causing new floaters or flashes. Sudden symptoms can occur with a retinal tear or detachment.

\medterm{Vitreous Hemorrhage} Bleeding into the vitreous cavity that causes sudden painless floaters, cobwebs, haze,
or loss of vision. Common causes include proliferative diabetic retinopathy, retinal
tear or detachment, trauma, and posterior vitreous separation.

\medterm{Vitreous Humor} The clear gel that fills the space between the lens and retina. It helps maintain the shape of the eye and transmits light to
the retina.

\medterm{VKORC1 Polymorphism} An inherited variation in the vitamin K epoxide reductase complex subunit 1 gene that can alter sensitivity to warfarin. Genotype may contribute to dose selection, but clinical factors, interacting medicines, diet, and INR monitoring remain important.
\medterm{VLDL Test} A blood test that measures or estimates very-low-density lipoprotein, a particle carrying triglycerides. The result may help assess cardiovascular risk but is interpreted with the full cholesterol and triglyceride profile.

\medterm{VO2 Max} The greatest amount of oxygen the body can use during intense exercise. It is a measure of aerobic fitness and is influenced by the heart, lungs, blood, muscles, age, health, and training.
\synonyms
Maximal Oxygen Consumption; Maximal Aerobic Capacity

\medterm{Vocal-Cord Dysfunction} Older term for inducible laryngeal obstruction, in which the larynx temporarily narrows during breathing and may cause throat tightness, noisy breathing, cough, or shortness of breath.

\medterm{Vocal Cord Nodules} Benign, usually paired thickened areas of the vocal folds caused by repeated voice strain or overuse. They can cause

hoarseness, vocal fatigue, and a rough or breathy voice.

\medterm{Vocal Cord Paralysis} Inability of one or both vocal folds to move normally, often because of nerve damage. It may cause hoarseness, weak cough, choking, or difficulty protecting the airway; sudden breathing trouble is an emergency.

\synonyms
Vocal Fold Paralysis

\medterm{Vocal Cord Paresis} Weakness or partial paralysis of one or both vocal cords, often from injury to the nerves that control them. It can cause a weak or breathy voice, voice fatigue, coughing during eating, or aspiration.

\medterm{Vocal Cord Polyp} A usually noncancerous growth on a vocal fold, often after heavy or sudden voice use. It can cause persistent hoarseness, a rough voice, or voice fatigue and is assessed by examining the larynx.

\medterm{Vocal Cords} Vocal cords, or vocal folds, are bands of tissue in the larynx that vibrate as air passes between them to produce the voice. Swelling, nodules, or injury can cause hoarseness.

\medterm{Vocal-Fold Nodules} Alternate name for vocal cord nodules, benign thickened areas of the vocal folds that can cause persistent hoarseness.

\medterm{Vocal Fold Paralysis} Alternate name; see \textit{Vocal cord paralysis}.

\medterm{Vocal-Fold Polyps} Alternate name for vocal cord polyps, usually noncancerous growths on the vocal folds that can cause persistent hoarseness or voice fatigue.

\medterm{Vocal Folds} Two folds inside the voice box that vibrate as air passes between them to make sound. They also close during swallowing to protect the lungs from food and liquid.

\medterm{Vocal Hemorrhage} Vocal hemorrhage is bleeding into a vocal fold, often after sudden or forceful voice use. It can cause abrupt hoarseness.

\medterm{Vocal Resonance} The sound of a person’s voice heard through the chest with a stethoscope while they speak. Changes in the sound can provide clues to pneumonia, fluid around the lung, a collapsed lung, or other lung disease.

\medterm{Vogt-Koyanagi-Harada Disease} A rare, multisystem, cell-mediated autoimmune disorder directed against melanocyte-associated antigens (such as tyrosinase), occurring predominantly in pigmented ethnic populations. Characterized by chronic bilateral granulomatous panuveitis with exudative retinal detachments, accompanied by sensorineural hearing loss, tinnitus, vitiligo, poliosis (whitening of eyelashes and hair), and aseptic meningitis.

\synonyms
VKH Syndrome; Uveomeningoencephalitic Syndrome

\medterm{Vogt-Koyanagi-Harada Syndrome} A multisystem autoimmune inflammatory disorder mediated by T-cell responses against melanocyte-associated antigens, characterized by bilateral chronic granulomatous panuveitis, exudative retinal detachments, meningismus (CSF pleocytosis), sensorineural dysacusis, poliosis, and patchy vitiligo.
\synonyms
VKH Syndrome; Uveomeningoencephalitic Syndrome

\medterm{Vohwinkel Syndrome} A rare inherited palmoplantar keratoderma causing thick skin on the palms and soles, constricting bands around fingers or toes, and sometimes hearing loss or other features. Severe constriction can threaten circulation and requires specialist care.

\medterm{Voice} Sound produced by vibration of the vocal folds as air passes through the larynx, then shaped by the throat, mouth, and nose for speech.

\medterm{Voice And Speech} The products of complex coordinated activity of the respiratory, laryngeal, and articulatory systems. Communication produced by coordinated breathing, voice-box vibration, and movements of the mouth, tongue, and other speech structures.

\medterm{Voice Box} The larynx, an organ in the neck containing the vocal folds and forming part of the airway. It produces voice and helps protect the lower respiratory tract during swallowing.

\medterm{Voice Disorders} Conditions affecting voice quality, pitch, loudness, or ease of speaking because of problems involving the vocal folds, airway, nerves, or voice use.

\synonyms
Dysphonia; Vocal Cord Dysfunction

\medterm{Void} To empty the urinary bladder by urinating; in other medical contexts, void may mean an empty space or absence of a substance.

\medterm{Voiding Cystourethrogram} An X-ray test that fills the bladder with contrast through a catheter and records the bladder and urethra while the person urinates. It can show urine flowing backward toward the kidneys, blockage, leakage, or structural problems, especially in children with selected urinary infections.

\medterm{Volar} Relating to the palm of the hand or sole of the foot, especially the surface opposite the dorsum.

\medterm{Volatile Anesthetics} Inhaled medicines used to produce and maintain unconsciousness during some operations. They can lower blood pressure and breathing and rarely cause a dangerous rise in body temperature in susceptible people.


\synonyms
Inhalational Anesthetics; Inhalation General Anesthetics

\medterm{Volkmann Contracture} A permanent, severe ischemic contracture of the forearm flexor musculature leading to a claw-like hand deformity, resulting from untreated acute compartment syndrome of the forearm following supracondylar humeral fractures, crush injuries, or tight casting.
\synonyms
Volkmann's Ischemic Contracture; Ischemic Myodystrophy

\synonyms Volkmann contracture; Volkmann's ischaemic contracture; Volkmann claw deformity.

\medterm{Volume Of Distribution} A calculated measure of how widely a medicine appears to spread through the body compared with the amount left in the blood. It helps guide dosing but does not represent a real space inside the body.

\medterm{Volume Overload} Excess total body sodium and water, causing peripheral or pulmonary oedema; common causes include heart failure, kidney failure, liver disease, and excessive intravenous fluid, and treatment depends on the cause and severity.

\medterm{Voluntary} Under conscious control or performed by choice, such as voluntary movement. Voluntary action is distinguished from reflex, autonomic, or involuntary activity.

\medterm{Voluntary Psychiatric Admission} Admission to a hospital or psychiatric facility with the person’s agreement for assessment, treatment, or safety support. Legal rules about consent, leaving, capacity, and conversion to involuntary status vary by jurisdiction.

\medterm{Volutrauma} Lung injury caused by overstretching the lungs, usually when a breathing machine delivers breaths that are too large. Careful ventilator settings can reduce this risk.

\medterm{Volvulus} Twisting of a loop of intestine that blocks passage through the bowel and may cut off its blood supply. It can cause severe pain, vomiting, swelling, and tissue death and can be life-threatening.

\medterm{Vomit} To eject stomach or upper gastrointestinal contents through the mouth; the expelled material may also be called vomitus. Persistent vomiting can cause dehydration and electrolyte abnormalities.

\medterm{Vomiting} Forceful emptying of stomach contents through the mouth. It may result from infection, digestive disease, medicines, pregnancy, migraine, brain disorders, or poisoning. Repeated vomiting can cause dehydration, disturbed body salts, or inhalation of vomit.

\medterm{Vomiting Blood} The passage of blood or coffee-ground material in vomit, usually indicating bleeding in the upper digestive tract. It can be a serious symptom.

\medterm{Vomitus} Material expelled from the stomach or upper gastrointestinal tract during vomiting. Its color and contents can provide clues but do not by themselves establish a diagnosis.

\medterm{Von Economo Encephalitis} An atypical form of epidemic encephalitis (encephalitis lethargica) characterized by pronounced lethargy, hypersomnolence, ophthalmoplegia, and a frequent late sequela of progressive post-encephalitic parkinsonism, pathologically marked by midbrain and basal ganglia inflammation.
\synonyms
Encephalitis Lethargica; Sleeping Sickness (Von Economo Type)

\medterm{Von Gierke Disease} Glycogen storage disease type I, an inherited deficiency of glucose-6-phosphatase or its transport system that causes fasting hypoglycemia, lactic acidosis, enlarged liver, and excess uric acid or lipids.

\medterm{Von Graefe Sign} A delay of the upper eyelid as the eyes look downward, leaving a strip of white visible above the iris. It can occur with an overactive thyroid and related eye disease.
\synonyms
Lid Lag Sign; Von Graefe's Lid Lag

\medterm{Von Hippel-Lindau Disease} An inherited condition, usually caused by a change in the \textit{VHL} gene, that increases the risk of tumors and cysts in the brain, spinal cord, eyes, kidneys, adrenal glands, and pancreas. These growths can affect vision, blood pressure, or kidney function.

\medterm{Von Hippel-Lindau Syndrome} Von Hippel-Lindau syndrome is an inherited disorder caused by changes in the VHL gene. It increases the risk of tumors and cysts in organs such as the brain, retina, kidneys, pancreas, and adrenal glands.

\medterm{Von Meyenburg Complex} A benign congenital biliary ductal plate malformation (biliary microhamartoma) consisting of small, non-progressive clusters of dilated, cystic bile ducts embedded in a dense collagenous stroma, usually discovered incidentally on liver imaging or biopsy.
\synonyms
Biliary Microhamartoma; Ductal Plate Malformation

\medterm{Von Recklinghausen Disease} Von Recklinghausen disease is another name for neurofibromatosis type 1, an inherited condition causing multiple café-au-lait spots, neurofibromas, and increased risk of certain nerve, bone, eye, and other problems.

\medterm{Von Recklinghausen’S Disease} Alternate name; see \textit{Neurofibromatosis type 1}.

\medterm{Von Recklinghausen’S Neurofibromatosis} Alternate name; see \textit{Neurofibromatosis type 1}.

\medterm{Von Recklinghausen’S Phakomatosis} Alternate name; see \textit{Neurofibromatosis type 1}.

\medterm{Von Rosen Splint} A malleable, lightweight orthopedic abduction splint applied in neonates with developmental dysplasia of the hip (DDH) to maintain the femoral head securely centered within the acetabulum in flexion and abduction.
\synonyms
Von Rosen Frame; DDH Abduction Splint

\medterm{Von Willebrand Disease} A bleeding disorder caused by too little or poorly working von Willebrand factor, a blood protein that helps clots form. It is often autosomal dominant, although severe type 3 and some other forms are autosomal recessive. It commonly causes nosebleeds, easy bruising, heavy menstrual bleeding, or prolonged bleeding after procedures.

\medterm{Voriconazole} An antifungal medicine used for selected serious fungal infections. It stops fungi making an essential part of their outer structure; vision changes, liver problems, medicine interactions, and skin sensitivity require monitoring.

\medterm{VOR Reflex} A reflex that moves the eyes opposite to head movement so vision stays steady while the head turns. Problems with this reflex can cause blurred or bouncing vision, dizziness, or imbalance.

\medterm{Voyeurism} A recurrent sexual interest in observing an unsuspecting person who is naked, disrobing, or engaging in sexual activity. The interest alone is not a mental disorder; voyeuristic disorder is the diagnostic term when required criteria are met.

\medterm{Voyeuristic Disorder} A paraphilic disorder involving recurrent, intense sexual arousal from observing an unsuspecting person who is naked, disrobing, or engaging in sexual activity. Diagnostic criteria include acting on the urges or clinically significant distress or impairment, with the required duration criteria.

\medterm{VPCs (Ventricular Premature Contractions)} Alternate name; see \textit{Premature ventricular contractions (PVCs)}.

\medterm{VRE (Vancomycin-Resistant Enterococci)} Enterococcus bacteria that are no longer killed by vancomycin, an antibiotic. They may live harmlessly in the body or cause urinary, wound, bloodstream, or other infections that require infection-control measures and suitable treatment.

\medterm{VSD (Ventricular Septal Defect)} A hole in the wall between the heart’s lower chambers. Small holes may close or cause no symptoms, while larger holes can make the heart and lungs work too hard and may need treatment.

\synonyms
Ventricular Septal Defect

\medterm{Vulnerary} A substance traditionally believed to help heal wounds. Such products may lack good evidence and can cause allergy or contamination.

\medterm{Vulva} The external female genital structures, including the mons pubis, labia, clitoris, and vestibule. It protects the openings of the urinary and reproductive tracts and contains specialized sensory and glandular tissue.

\synonyms
Female External Genitalia; Pudendum

\medterm{Vulval Intraepithelial Neoplasia} Vulval intraepithelial neoplasia is abnormal cell growth limited to the surface layers of the vulvar skin. Some types are linked to HPV, and untreated lesions may progress to invasive cancer.

\medterm{Vulvar} Relating to the vulva, the external female genital structures including the labia, clitoris, vestibule, and associated glands.

\medterm{Vulvar Cancer} Cancer arising in the external female genitalia. It may appear as a persistent lump, ulcer, wart-like lesion, itching, pain, or bleeding.

\medterm{Vulvar Carcinoma} Cancer of the external female genital area, most often arising from surface skin cells. It may appear as a persistent lump, sore, wart-like growth, itching, pain, or bleeding.

\synonyms
Vulvar Cancer; Squamous Cell Carcinoma of Vulva

\medterm{Vulvar Inclusion Cysts} Small sacs under the skin of the vulva that contain tissue trapped from the surface lining. They may develop after childbirth tears, surgery, or other injury, or occur without an obvious cause. Most cause no symptoms, but they may become swollen, painful, infected, or uncomfortable during sexual activity. Symptomatic cysts may need removal.

\medterm{Vulvar Intraepithelial Neoplasia} Abnormal cells limited to the surface layer of the vulva. Some forms can progress to invasive cancer, especially without treatment. It may cause itching, burning, sores, warty growths, or a color change, but some people have no symptoms.

\medterm{Vulvar Melanoma} A rare cancer of the pigment-producing cells in the vulva. It may appear as a dark or skin-colored lump, patch, sore, or bleeding area; diagnosis requires a biopsy, and treatment depends on size and spread.

\medterm{Vulvar Squamous Cell Carcinoma} The most common invasive vulvar malignancy, arising through HPV-related or HPV-independent pathways. It usually presents as a persistent mass, ulcer, plaque, or wart-like lesion and requires biopsy for diagnosis.

\medterm{Vulvectomy} Surgical removal of part or all of the vulva, performed for selected cancers, precancerous lesions, or severe persistent disease.

\medterm{Vulvitis} Inflammation or irritation of the external female genital area. It can cause itching, burning, redness, swelling, soreness, or discharge and may result from infection, allergy, skin disease, hormones, or friction.

\medterm{Vulvodynia} Vulvodynia is chronic vulvar pain without an identifiable specific cause, often described as burning or rawness and triggered by touch, pressure, or occurring spontaneously.

\synonyms
Chronic Vulvar Pain

\medterm{Vulvovaginal Candidiasis} A yeast infection of the vulva and vagina, usually caused by Candida. It commonly causes intense itching, burning, soreness, redness, and thick discharge; recurrent, severe, unusual, or treatment-resistant symptoms need testing to confirm the cause.

\medterm{Vulvovaginitis} Simultaneous inflammation or infection of the vulva and vagina, commonly presenting with pruritus, burning, dysuria, dyspareunia, and abnormal vaginal discharge, resulting from infectious etiologies (such as Candida albicans, Bacterial vaginosis, Trichomonas vaginalis) or non-infectious causes (atrophic vaginitis, contact irritants).
\synonyms
Vaginitis and Vulvitis; Vulvovaginal Inflammation
